\documentclass[11pt,a4paper]{article}
\usepackage[margin=27mm]{geometry}
\usepackage[T1]{fontenc}
\usepackage{lmodern,microtype,amsmath,amssymb,amsthm,mathtools,mathrsfs}
\usepackage{booktabs,longtable,array,enumitem,xcolor}
\usepackage[authoryear,round]{natbib}
\usepackage{hyperref}
\hypersetup{colorlinks=true,linkcolor=blue!45!black,citecolor=blue!45!black,urlcolor=blue!45!black}
\setlength{\emergencystretch}{2em}
\setlist{nosep}
\newcommand{\statement}[3]{\par\bigskip\noindent\textbf{#1. [#2] #3.}\quad\phantomsection\label{stmt:#1}}
\newcommand{\Part}[2]{\section*{Part #1. #2}\addcontentsline{toc}{section}{Part #1. #2}}
\newcommand{\wh}[1]{\widehat{#1}}
\newcommand{\ol}[1]{\overline{#1}}
\newcommand{\GT}{\widehat{\mathrm{GT}}}
\newcommand{\GTzero}{\widehat{\mathrm{GT}}_0}
\newcommand{\cP}{\mathcal P}
\newcommand{\cH}{\mathcal H}
\newcommand{\cK}{\mathcal K}
\newcommand{\Z}{\mathbb Z}
\newcommand{\Q}{\mathbb Q}
\newcommand{\F}{\mathbb F}
\newcommand{\C}{\mathbb C}
\newcommand{\R}{\mathbb R}
\DeclareMathOperator{\Aut}{Aut}
\DeclareMathOperator{\Out}{Out}
\DeclareMathOperator{\Inn}{Inn}
\DeclareMathOperator{\im}{im}
\DeclareMathOperator{\Hom}{Hom}
\DeclareMathOperator{\Gal}{Gal}
\DeclareMathOperator{\SL}{SL}
\DeclareMathOperator{\PSL}{PSL}
\DeclareMathOperator{\Sym}{Sym}
\DeclareMathOperator{\rank}{rank}
\DeclareMathOperator{\cl}{cl}
\title{LM4\_SUCCESSOR\\Involutions, the icosahedral centralizer,\\and the remaining comparison problems}
\author{AI-produced by GPT-6 Astra (OpenAI, ChatGPT Pro mode)}
\date{16 November 2026\\\textbf{INCOMPLETE: research in progress}}
\begin{document}
\maketitle
\begin{abstract}
We audit and develop the involution and finite-centralizer problems posed in the accompanying LM4 successor prompt. All conclusions are assigned explicit proof statuses. No priority is claimed. This working version is incomplete: the unresolved questions remain explicit, and a failure of a proposed argument is never treated as a counterexample to its conclusion.
\end{abstract}
\tableofcontents
\clearpage
\section*{Scope and conventions}
We distinguish full profinite completion, pro-$p$ completion, relative pro-$p$ completion, and procongruence completion throughout. A subgroup closure is taken in the explicitly named ambient completion. Continuous homomorphisms and nonabelian continuous cohomology are understood unless a discrete category is specified. Earlier AI-produced returns are cited with their run names; the supplied reports are not substitutes for proofs. The specified refereed inputs LM4 D.32--D.34 and LM7\_PUSH D.50, E.58 and E.65 are identified wherever used. The Montpellier transcriptions are provisional readings, not the manuscript scans.
\Part{A}{Involutions of the three-cusp group}
\section{Involutions of the three-cusp group}

Throughout this part \(F=\widehat F(x,y)\), \(z=(xy)^{-1}\), and
\(\mathcal P\) denotes the subgroup of \(\operatorname{Out}(F)\)
preserving each named peripheral conjugacy class up to a common
unit power. All group homomorphisms are continuous. We use the
normalized lift
\[
 (\lambda,f):\quad x\longmapsto x^\lambda,\qquad
 y\longmapsto f^{-1}y^\lambda f,\qquad f\in F',
\]
where \(F'\) is the closed derived subgroup. This normalization is
a homomorphic section, as in the refereed LM4 D.32. Write
\(c=(-1,1)\). We use R1--R3 only with their stated scopes.

\statement{A.1}{Known}{Torsion in an acting group and its action on torsion are different assertions}
Lochak and Schneps (1997, Proposition 4(i)) prove
\(C_F(c)=1\); Proposition 4(ii) proves that \(c\) is
self-centralizing in \(\widehat{GT}\).
Schneps (1997, \S1.4(4)) and Lochak and Schneps (2006,
Question 1.4) ask whether all nonidentity torsion in
\(\widehat{GT}\) consists of conjugates of \(c\).
Harbater and Schneps (2000, Main Theorem(a), proved by
Propositions 3--4 and Theorem 2) identify
\(\widehat{GT}_0\) with the centralizer in \(\mathcal P\) of
the geometric \(S_3\).
Collas's genus-zero Theorems 3.3 and 3.10 identify prime-order
torsion in \(\widehat\Gamma_{0,[n]}\) with discrete torsion
and prove the \(\lambda\)-conjugacy action of \(\widehat{GT}\).
His genus-one Theorems 6.3 and 6.5 give the corresponding
statements with acting group \(\widehat{GR}\), isomorphic
to the subgroup denoted \(I\Gamma\) in that article;
replacing it there by all of \(\widehat{GT}\) is not
justified by those theorems.\cite{A_LS1997,A_Schneps1997,A_LS2006,A_HS2000,A_Collas2012a,A_Collas2012b}

\begin{proof}
For the centralizer comparison identify \(K(0,4)\) with
\(\widehat F(x,y)\). The natural \(S_4\)-action on its outer
automorphism group factors through the anharmonic \(S_3\).
Harbater--Schneps's quasi-special condition is precisely the
three peripheral power condition here; their symmetric
condition is commutation with that \(S_3\). Their proof
translates commutation with the transposition into the
two-cycle relation and commutation with the three-cycle into
the hexagon. The identification has the required source
and target, not merely an abstract isomorphism.

In Collas's formulas the variable of prime order lies in a
completed mapping class group; \((\lambda,f)\) is an arbitrary
element of the acting group. The conclusion
\(F(\gamma)=u\gamma^\lambda u^{-1}\) does not classify
finite-order elements \(F\) of the acting group. The subgroup
\(\widehat{GR}\) in the genus-one article is retained in its
Theorem 6.5. Neither action theorem therefore proves task A.8.
\end{proof}

The opened sources establish no literature antecedent for the
full order bound in R1. Searches included the authors' source
lists, the cited torsion papers, and recent work on profinite
curve-complex automorphisms and Galois--Teichm\"uller theory.
This records the search, not a priority claim or a proof that
no antecedent exists.

\statement{A.2}{Proved}{Relative pronilpotent levels recover the full peripheral group}
For an open characteristic subgroup \(U\) of \(F\), put
\[
 R_U=R_{\mathrm{nil}}(U),\quad F_U=F/R_U,\quad
 P_U=\operatorname{Out}^{\sharp}(F_U).
\]
For \(V\subset U\), the quotient \(F_V\to F_U\) induces
a continuous epimorphism \(P_V\to P_U\). There is an
isomorphism of profinite groups
\[
                 \mathcal P \xrightarrow{\sim}\varprojlim_U P_U.
\]
Two elements of \(\mathcal P\) are conjugate in \(\mathcal P\)
if and only if their images are conjugate in every \(P_U\).
This proves Lead 1.

\begin{proof}
Restricting finite nilpotent quotients of \(U\) to \(V\)
shows \(R_V\subset R_U\). Each \(R_U\) is characteristic in
\(F\). Since \(R_U\subset U\) and open characteristic
subgroups are cofinal in the finitely generated profinite
group \(F\), we have \(\bigcap_U R_U=1\).

We first prove surjectivity onto \(P_U\). Choose a representative
of a peripheral outer class with common multiplier \(\lambda\),
and a Galois element with cyclotomic multiplier \(\lambda\);
the cyclotomic character is onto
\(\widehat{\mathbb Z}^{\times}\). The latter acts on \(F_U\).
Composing with its inverse gives a multiplier-one peripheral
automorphism. R2 lifts this automorphism, with chosen
conjugators, to \(F\). Composing back proves surjectivity.
In particular every peripheral automorphism of \(F_V\)
lifts to \(F\), so preserves the image of \(R_U\).
This proves that the transition maps are defined on the
entire displayed outer groups. Their surjectivity follows
by lifting first to \(\mathcal P\).

Here the automorphism groups are compact: intersecting
kernels of all maps to finite groups of bounded order
gives a cofinal system of open fully invariant subgroups.
The conditions that a continuous endomorphism induces
bijections on these finite quotients are closed conditions
on its generator images. Consequently \(\operatorname{Aut}(F)\)
is profinite. Inner automorphisms form a compact subgroup.
The peripheral conditions are closed, since conjugators
and the common multiplier range over compact groups.
The same argument applies to the finitely generated
quotients \(F_U\).

If \([a]\in\mathcal P\) dies in every \(P_U\), the sets
\[
 \{g\in F:a(x)\equiv gxg^{-1},\
              a(y)\equiv gyg^{-1}\pmod {R_U}\}
\]
are nonempty compact sets with the finite intersection
property. Their intersection gives an actual inner
representative for \(a\), proving injectivity.
For a compatible family of classes in \(P_U\), their
nonempty compact inverse images in \(\mathcal P\) have
the same finite intersection property, proving surjectivity.
Finally apply this argument to the transporter sets
\(\{g:g_Ua_Ug_U^{-1}=b_U\}\) in \(\mathcal P\).
Surjectivity onto \(P_U\) makes these nonempty whenever
the two images are conjugate. Their intersection is the
required full conjugator.
\end{proof}

\statement{A.3}{Proved}{A central quotient is the exact obstruction to a peripheral lift}
Let \(N\) be a closed characteristic subgroup of \(F\), and
let \([N,F]\) denote the closed subgroup generated by
commutators between \(N\) and \(F\).
An automorphism \(a\) of \(F/N\), with specified conjugators
\(b_i\) satisfying
\[
 a(\bar x_i)=b_i\bar x_i b_i^{-1},\qquad
 (x_1,x_2,x_3)=(x,y,z),
\]
lifts to a peripheral automorphism of \(F\) with these
conjugators if and only if it lifts to a peripheral
\emph{automorphism} of \(F/[N,F]\) with conjugators above
the \(b_i\). This proves Lead 2.

\begin{proof}
Necessity follows from characteristicity. For sufficiency
let \(a^\#\) and \(b_i^\#\) be the data on \(F/[N,F]\).
We use the finite central-refinement count in the refereed
LM7\_PUSH D.50, equations (D.50.6)--(D.50.16).
For any actual quotient \(q_E:F\twoheadrightarrow E\)
and normal subgroup \(A\triangleleft E\), that count
constructs an actual quotient
\(q^\dagger:F\twoheadrightarrow D^\dagger\), with central
kernel over \(D=E/A\). A second epimorphism to
\(D^\dagger\) with specified peripheral conjugators has
the same number of epimorphism lifts to \(E\), with the
transported \(A\)-conjugacy classes on the boundary.
This counting assertion uses only normality of \(A\).
The hypothesis \([E,A]=A\) occurs in the separate
joint-image assertion (D.50.17), not used here.

Take \(A=q_E(N)\). The image \(q^\dagger(N)\) is central,
so \(q^\dagger\) kills \([N,F]\).
Thus \(q^\dagger a^\#\) is an epimorphism from the
same source to \(D^\dagger\), with conjugators the
images of \(b_i^\#\). The finite count gives an epimorphism
to \(E\) whose conjugators agree with the \(b_i\) in \(D\);
the reference count is positive since it includes \(q_E\).

In the compact space \(F^3\) of conjugators impose the
boundary product relation, generation in every finite
quotient, and agreement with each \(b_i\) in every finite
quotient of \(F/N\). A finite list of conditions lies in
one joint finite quotient \(E\), and the preceding
construction satisfies that list. Compactness gives
full conjugators satisfying all conditions. Their
endomorphism is surjective since its compact image
maps onto every finite quotient; Hopficity makes it
an automorphism. Its generator images induce exactly
\(a\) modulo \(N\).
\end{proof}

\statement{A.4}{Proved}{A degree-three coefficient imposed by the third cusp}
Over a field \(k\) of characteristic zero or greater
than three, use the free class-three Lie algebra on
\(X,Y\), and set
\[
 Z=[X,Y],\quad U_3=[X,Z],\quad V_3=[Y,Z],\quad
 x=\exp X,\quad y=\exp Y .
\]
If \(x\mapsto x^\mu,\ y\mapsto g^{-1}y^\mu g\), with
\(\mu\in k^\times\) and
\(\log g=tZ+\text{central terms}\), sends \(xy\)
to a conjugate of \((xy)^\mu\), then
\(t=(\mu^2-1)/24\).

\begin{proof}
The Baker--Campbell--Hausdorff formula gives
\[
 W=\log(xy)=X+Y+\tfrac12Z+\tfrac1{12}(U_3-V_3).
\]
The logarithm of the image is
\[
 \mu(X+Y)+\tfrac{\mu^2}{2}Z+
       \tfrac{\mu^3}{12}(U_3-V_3)+\mu tV_3 .
\]
Write the logarithm of a conjugator as
\(aX+bY+dZ+\text{central terms}\).
In \(\exp(\operatorname{ad}(aX+bY+dZ))(\mu W)\),
comparison of degree two gives \(m=a-b=(\mu-1)/2\).
The difference of the \(U_3\)- and \(V_3\)-coefficients,
divided by \(\mu\), then gives
\[
 \frac{\mu^2}{6}-t=\frac16+\frac m2+\frac{m^2}{2}
                  =\frac{3\mu^2+1}{24}.
\]
The result follows. The coefficient \(d\) and all
central coordinates cancel.
\end{proof}

\statement{A.5}{Refuted}{The specified conjugator in task A.3 need not preserve the third cusp}
There is an involution \(\alpha\in\widehat{GT}\), conjugate
to \(c\) by a Galois element, for which the unique solution
\(h\in F'\) of \(f_\alpha=h^{-1}\alpha(h)\) gives
\[
                 \psi(x)=x,\qquad\psi(y)=hyh^{-1}
\]
that does not preserve the conjugacy class of \(z\)
on \(F/R_{\mathrm{nil}}(F)\). The failure occurs in
the class-three exponent-seven quotient of order \(7^5\).

\begin{proof}
Choose \(\theta=(\mu,g)\) in the Galois image with
\(\mu\equiv2\pmod7\), possible by surjectivity of the
cyclotomic character modulo seven, and put
\(\alpha=\theta c\theta^{-1}\).
Normalization is a homomorphic section, so these
are actual automorphisms with that equality.
From \(\alpha\theta(y)=\theta(y)^{-1}\) we obtain
\[
 \alpha(y)=\alpha(g)g^{-1}y^{-1}g\alpha(g)^{-1},
 \qquad f_\alpha=g\alpha(g)^{-1}.
\]
Thus \(h=g^{-1}\) is a solution. Also
\(C_F(\alpha)=\theta(C_F(c))=1\) by
Lochak--Schneps Proposition 4(i). The quotient of two
solutions is fixed by \(\alpha\), so \(h\) is unique.

Let \(Q=\exp L\) for the free class-three Lie algebra
\(L\) on two generators over \(\mathbb F_7\).
It has order \(7^5\), exponent seven, and is generated
by \(x=\exp X,y=\exp Y\); hence it is a quotient of
\(F^{(7)}\) and of \(F/R_{\mathrm{nil}}(F)\).
The third-cusp condition on \(\theta\) and A.4 give
\[
       \operatorname{coeff}_Z(\log\bar g)
                    =\frac{2^2-1}{24}=1 .
\]
But \(\psi\) has multiplier one and the same conjugator
\(g\). Peripheral preservation for \(\psi\) would force
this coefficient to be zero by A.4, a contradiction.

The exact certificate \texttt{check\_A3\_class3.py}
realizes \(Q\) in the units of
\[
 \mathbb F_7\langle X,Y\rangle/(\text{words of length }4).
\]
It enumerates all \(16807\) elements; the conjugacy
orbit of \(xy\) has \(49\) elements. For
\(g=\exp[X,Y]\) it verifies
\[
 xg^{-1}yg\not\sim xy,\qquad
 x^2g^{-1}y^2g\sim(xy)^2 .
\]
The actual Galois conjugator differs only by central
degree-three coordinates relevant to this calculation,
which do not affect it. This certifies a failure in
an actual full profinite example; no order-preserving
lift of a finite involution is being assumed.
\end{proof}

The correct universal statement from R1 is that this
\(\psi\) intertwines \(\alpha,c\) on the maximal
pro-\(p\) quotients as automorphisms of free groups.
Peripheral preservation of this specific \(\psi\)
is an additional condition and is false.

\statement{A.6}{Proved}{The pure centralizer of standard conjugation has two elements}
\[
 C_{\mathcal P}(c)=N_{\mathcal P}(\langle c\rangle)=\{1,c\}.
\]
If \(\alpha\in\widehat{GT}_0\), every conjugator in
\(\mathcal P\) from \(c\) to \(\alpha\) belongs
to \(\widehat{GT}_0\).

\begin{proof}
A normalized \((\mu,g)\) commuting with \(c\) satisfies
\(c(g)=g\), by the normalized group law. A.1 gives
\(g=1\). In particular the multiplier-one centralizer
is trivial. Let \(\Sigma=S_3\) be the geometric
permutation group. If
\(\beta c\beta^{-1}=\alpha\) and \(\alpha,c\) commute
with \(\Sigma\), then
\(r_\sigma=\beta^{-1}\sigma\beta\sigma^{-1}\)
is pure, has multiplier one, and centralizes \(c\).
Thus \(r_\sigma=1\), and the Harbater--Schneps
identification in A.1 puts \(\beta\) in
\(\widehat{GT}_0\).

Apply this to a centralizing \((\mu,1)\).
The hexagon, with \(m=(\mu-1)/2\), reads
\((xy)^m=y^mx^m\). In the finite Heisenberg groups
with multiplication
\((a,b,d)(a',b',d')=(a+a',b+b',d+d'+ab')\),
the final coordinates of these two words are
\(m(m+1)/2\) and zero. Hence \(\mu^2=1\) in
\(\widehat{\mathbb Z}\). The actual automorphism
\((\mu,1)\) then has square one. R1 forces its
multiplier to be globally \(1\) or \(-1\), proving
the centralizer assertion. A normalizer of
\(\langle c\rangle\) centralizes its unique
nonidentity element.
\end{proof}

\statement{A.7}{Proved}{The multiplier-one cohomology set already has continuum many classes}
For \(G=\mathcal P,\widehat{GT}_0\), or \(\widehat{GT}\),
put \(H_G=\ker\lambda\), with \(C_2\) acting by \(c\).
There is a canonical closed injection of profinite spaces
\[
 \widehat{\mathbb Z}^{\times}/\{1,-1\}
 \longrightarrow H^1(C_2,H_G),\qquad
 \mu\longmapsto[\theta_\mu c\theta_\mu^{-1}c],
\]
where \(\theta_\mu\) is any element of multiplier \(\mu\).
The set \(H^1(C_2,H_G)\) has cardinality \(2^{\aleph_0}\).
For \(G=\widehat{GT}_0\), \(\Delta\) vanishes on this
closed family. All these \(H_G\)-conjugacy classes form
one \(G\)-conjugacy class.

\begin{proof}
A cocycle is \(a\in H_G\) with \(a\,c(a)=1\); it
corresponds to the involution \(ac\). Cohomology
identifies \(a\) and \(h^{-1}a\,c(h)\), exactly
conjugacy by \(H_G\). The Galois image makes
\(\lambda:G\to\widehat{\mathbb Z}^{\times}\) onto.
The element \(a_\mu=\theta_\mu c\theta_\mu^{-1}c\)
is a cocycle, and replacing \(\theta_\mu\) by a
different lift conjugates its involution by \(H_G\).

Two displayed involutions are \(H_G\)-conjugate
if and only if
\(\theta_\nu^{-1}h\theta_\mu\in C_G(c)=\{1,c\}\)
for some \(h\in H_G\), by A.6. Taking multipliers
gives \(\mu=\pm\nu\); these conditions are sufficient
as well. This proves the injection.
The cocycle space is compact, and its quotient by
the compact conjugating group is Hausdorff.
The orbit map from \(G\) factors continuously
through the multiplier quotient and then its
two signs, so the injection is continuous and closed.
Its domain has cardinality continuum. Conversely
\(G\) is contained in the metrizable profinite group
\(\operatorname{Out}(F)\), of cardinality at most
continuum. This proves equality.
Finally every represented involution is conjugate
to \(c\); the character \(\Delta\) is conjugacy
invariant and \(\Delta(c)=0\).
\end{proof}

\statement{A.8}{Proved}{Full conjugacy is an additional quotient of nonabelian cohomology}
With the notation of A.7, the full \(G\)-conjugacy
classes of nonidentity involutions are
\[
 \widehat{\mathbb Z}^{\times}\backslash H^1(C_2,H_G).
\]
The action of \(\mu\), using any lift \(\theta_\mu\), is
\[
 [a]\longmapsto[\,\theta_\mu a\,c(\theta_\mu^{-1})\,].
\]
For \(\mathcal N=\ker(\mathcal P\to P_F)\), the
\(\mathcal N\)-conjugacy classes of involutions with
image exactly \(c_F\) are \(H^1(C_2,\mathcal N)\).
Full \(\mathcal P\)-conjugacy on that fibre further
quotients by \(C_{P_F}(c_F)\), acting by lifted
conjugation, which need not preserve the distinguished
cohomology class.

\begin{proof}
By R1 every nonidentity involution has multiplier
\(-1\), so is uniquely \(ac\), \(a\in H_G\).
Its square-one condition is the cocycle equation.
Conjugation by \(\theta_\mu\) gives the formula.
Different lifts differ by \(H_G\), already divided
out in cohomology, proving the first assertion.

The same calculation for \(\mathcal P/\mathcal N\)
shows that involutions over the literal \(c_F\)
are \(ac\), \(a\in\mathcal N\). A conjugation
preserves this fibre exactly when its quotient
centralizes \(c_F\). Changing a lift changes
the action only by \(\mathcal N\)-conjugacy.
Such a lift need not centralize \(c\), explaining
the last qualification.
\end{proof}

\statement{A.9}{Proved}{Odd-prime conjugacy over a relative Frattini level}
Let \(U\) be open characteristic in \(F\), let \(p\)
be odd, put \(E_p=F/O^p(U)\), \(V=U/O^p(U)\),
and consider the peripheral outer subgroup preserving
\(V\). Its kernel on
\[
           E_p/\Phi(V)=F/\overline{U'U^p}
\]
is pro-\(p\). Two involutions with the same image
in this Frattini quotient are conjugate by the kernel.
This applies to the fibre over \(c\), where that
fibre is defined by descent to \(F/U\).

\begin{proof}
\(V\) is a nonabelian finitely generated free
pro-\(p\) group. Let \(J\) be the actual
automorphisms of \(E_p\) inducing identity on
\(E_p/\Phi(V)\). Restriction maps \(J\) into
the Frattini kernel of \(\operatorname{Aut}(V)\).
That kernel is pro-\(p\): on finite lower
\(p\)-central quotients its successive kernels
embed into homomorphism groups of elementary
abelian \(p\)-groups, and those quotients are cofinal.

Restriction is injective. If an element of \(J\)
fixes \(V\) pointwise, write \(a(g)=gu_g\) with
\(u_g\in\Phi(V)\). Comparing the conjugation
actions on \(V\) gives \(u_g\in C_V(V)=1\).
Thus \(J\) is pro-\(p\). If an outer automorphism
acts inner on \(E_p/\Phi(V)\), lift its inner
conjugator to \(E_p\) and correct it to an
element of \(J\). Consequently the outer
kernel is a continuous quotient of a closed
subgroup of \(J\), hence pro-\(p\).

The common image of the two involutions has
order two, since a pro-\(p\) group for odd \(p\)
has no involution. In its inverse image the
two cyclic subgroups are complements to the
normal pro-\(p\) kernel. Finite
Schur--Zassenhaus gives conjugators in every
finite quotient; compactness gives a conjugator
in the kernel itself.
\end{proof}

For arbitrary quotient automorphisms, preservation
of the image of a merely characteristic \(U\)
must be checked. It is included in the phrase
over the class of \(c\). For fully characteristic
\(U\) it is automatic: any quotient endomorphism
lifts to an endomorphism of the free profinite
group and preserves \(U\). Such \(U\) are cofinal.

\statement{A.10}{Proved}{Every negative pure pro-\(p\) involution is locally conjugate for odd \(p\)}
Let \(P^{(p)}\) be the pure peripheral outer group
of \(F^{(p)}\), with multiplier in
\(\mathbb Z_p^\times\), for odd \(p\).
Every involution of multiplier \(-1\) is
conjugate to \(c\) by a multiplier-one element.

\begin{proof}
The multiplier kernel \(H^{(p)}\) acts trivially
on abelianization, so lies in the pro-\(p\)
Frattini kernel used in A.9. The inverse image
of \(\{1,-1\}\) under the multiplier is
\(H^{(p)}\rtimes\langle c\rangle\).
Every negative involution defines a complement;
Schur--Zassenhaus gives the required conjugacy.
\end{proof}

This does not rescue the particular map of A.5.
Its local coboundary can be changed by a local
fixed element, whereas its full coboundary is
unique. The comparison of fixed subgroups
need not be surjective.

\statement{A.11}{Proved}{A peripheral reflection triangle is equivalent to full conjugacy}
For a nontrivial involution \(\alpha\in\mathcal P\)
let \(E=F\rtimes_\alpha C_2\). Then \(\alpha\)
is conjugate to \(c\) in \(\mathcal P\) exactly
when there are involutions \(t_0,t_1,t_2\)
representing its three section classes such that
\[
 a=t_1t_0,\qquad b=t_0t_2,\qquad d=t_2t_1
\]
generate conjugates of the respective inertia
groups of \(x,y,z\), and the \(t_i\) generate
\(E\) topologically. The last condition is
equivalent to generation in every \(E/R_U\).

\begin{proof}
For standard \(c\) take
\(t_0=c,t_1=xc,t_2=y^{-1}c\); their products
are \(x,y,z\) and they generate the extension.
A peripheral conjugator transports this triple.

Conversely the universal property gives a
surjection from the profinite coproduct
\(C_2^{(0)}\amalg C_2^{(1)}\amalg C_2^{(2)}\)
to \(E\), sending the factors to the named
involutions. Its parity kernel is free
profinite on \(a,b\). The restriction maps
onto \(F\): every even word in three
reflections is a word in the two edge products.
Both kernels are free profinite of rank two,
so Hopficity makes this restriction an
isomorphism. The equality \(abd=1\) is literal.
Abelianization makes the unit exponents in
the three peripheral conditions equal,
so the induced outer map belongs to \(\mathcal P\).
Conjugation by \(t_0\) inverts both \(a\) and \(b\).
Its outer action on \(F\) is \(\alpha\), since
\(t_0\) and the named splitting element have the
same image in \(E/F\). Hence the resulting outer
map conjugates \(c\) to \(\alpha\); no choice of
an actual splitting representative is needed.

Finally put \(H=\overline{\langle t_0,t_1,t_2\rangle}\).
If \(HR_U=E\) for all \(U\), then for every \(e\in E\)
the compact sets \(H\cap eR_U\) are nonempty
and have the finite intersection property.
Their intersection and \(\bigcap R_U=1\)
imply \(e\in H\). Thus \(H=E\).
\end{proof}

The transcription of cote 145, scan 23, lists
\(l_0=t_\infty t_1,\ l_1=t_0t_\infty,\
l_\infty=t_1t_0\). With ordinary multiplication
these satisfy \(l_0l_\infty l_1=1\).
The cyclic order or path-composition convention
must therefore be made explicit before
identifying them with \(x,y,z\).
We use \(abd=1\), and do not emend an unseen scan.

\statement{A.12}{Open}{Dyadic and mixed-prime involution obstructions}
A.9 does not determine all two-primary classes or all
gluing classes over \(U=F'F^3\) and \(U_{12}\).
A.25 below constructs an actual involution with
parity one in the image of \(\widehat{GT}_0\) at
\(U_{12}\), and A.24 proves that this parity
descends to that outer image. Thus the existence
assertion at this specified relative level is proved.
Lead 3 still needs a compatible order-two full lift.
Neither R2 nor A.3 supplies that further equivariance.

For a surjection \(G\to Q\) with kernel \(L\),
a chosen lift \(t\) of an involution can be
corrected to an involution \(ut\), \(u\in L\),
exactly when
\[
                    u(tut^{-1})=t^{-2}.
\]
Surjectivity asserts the existence of \(t\),
not solvability of this nonabelian norm equation.
The map \(C_4\to C_2\) disproves that implication
in general, but is not a counterexample in
the peripheral groups. The residual task is
to solve the norm equations compatibly in
the actual relative-level kernels with the
specified frames and \(\Delta=1\), or to
prove they are insoluble.

\statement{A.13}{Open}{Parity on arbitrary involutions}
For an involution \(\tau=(-1,f)\) in
\(\widehat{GT}_0\), R3 gives
\(\Delta(\tau)=\delta((1,f))\).
It is zero for A.7's continuum family and
for every \(\widehat{GT}_0\)-conjugate of \(c\).
The cocycle relation alone gives only
\(2\Delta(a)=0\) in \(\mathbb F_2\).
A.25 gives such an involution in an actual relative
outer image of \(\widehat{GT}_0\), but no full
involution of parity one or universal full vanishing
proof is obtained, in \(\widehat{GT}_0\) or
\(\widehat{GT}\).

\statement{A.14}{Open}{The three full conjugacy assertions}
We have not decided whether every involution
of \(\mathcal P,\widehat{GT}_0,\widehat{GT}\)
is conjugate to \(c\) in its respective group.
A.2 detects conjugacy of a fixed pair on
all relative levels; A.11 identifies the
equivalent peripheral generation condition.
Neither supplies the compatible local
conjugators or reflection triangles whose
existence is required.

For \S31, folios 198--199 verso, scans
391--394, the full finite-subgroup conclusion
proved by R1 is that the only abstract
possibilities in \(\mathcal P\) are \(1,C_2\),
and \(C_2=\langle c\rangle\) occurs.
Exhaustion of its embedding conjugacy
classes is exactly the remaining assertion.

\statement{A.15}{Proved}{The pronilpotent lifting kernel also has continuum many cohomology classes}
Put \(Q=\mathcal P_F\), \(\mathcal N=\ker(\mathcal P\to Q)\),
and write \(c_Q\) for the image of \(c\). There is a canonical
closed injection
\[
 C_Q(c_Q)/\langle c_Q\rangle
 \longrightarrow H^1(C_2,\mathcal N),\qquad
 \bar\beta\longmapsto[\beta c\beta^{-1}c].
\]
Its image is exactly the \(\mathcal N\)-conjugacy classes in the
fibre over \(c_Q\) which are conjugate to \(c\) in \(\mathcal P\).
Both its domain and \(H^1(C_2,\mathcal N)\) have cardinality
\(2^{\aleph_0}\). Thus the large cohomology set does not itself
provide an additional full involution-conjugacy class.

\begin{proof}
The pronilpotent quotient is the product of its maximal
pro-\(p\) quotients. Every factor is its characteristic
Sylow subgroup; conversely, a tuple of continuous
automorphisms of the factors gives a continuous automorphism
of their product. Inner automorphisms and the three peripheral
conditions are coordinatewise. Therefore
\[
 Q=\prod_p\mathcal P^{(p)},\qquad
 C_Q(c_Q)=\prod_p C_{\mathcal P^{(p)}}(c).
\]
For odd \(p\), the multiplier map on the latter centralizer
is onto \(\mathbb Z_p^\times\). Indeed choose an element
\(\theta\) of the desired multiplier using the Galois image.
By A.10, some multiplier-one \(v\) conjugates
\(\theta c\theta^{-1}\) to \(c\); then \(v\theta\)
centralizes \(c\) and retains that multiplier.
Consequently the displayed centralizer, and its quotient
by the finite group \(\langle c_Q\rangle\), have at least
continuum many elements.

By A.2 every \(\bar\beta\in C_Q(c_Q)\) lifts to
\(\beta\in\mathcal P\). Its displayed commutator with
\(c\) lies in \(\mathcal N\) and is a cocycle.
Changing the lift changes its class only by
\(\mathcal N\)-conjugacy. Two lifts \(\beta_1,\beta_2\)
give the same class exactly when
\(\beta_2^{-1}n\beta_1\in C_{\mathcal P}(c)=\{1,c\}\)
for some \(n\in\mathcal N\). On taking the quotient,
this is exactly equality of the two cosets modulo
\(\langle c_Q\rangle\); the converse follows by lifting
that equality. This proves injectivity and the asserted
description of the image.
Compactness makes the injection continuous and closed,
as in A.7. All groups and orbit spaces here are compact
metrizable, so their cardinalities are at most continuum.
The preceding lower bound proves the asserted equalities.
\end{proof}

If the same construction is made in \(\widehat{GT}_0\),
using its actual image in any \(c\)-stable quotient,
every class in the connecting-map image has \(\Delta=0\).
This follows from conjugacy invariance of the character;
no surjectivity onto the entire peripheral quotient is
required for this observation.

\statement{A.16}{Proved}{A prescribed reflection splits every pro-two extension of the symmetric group}
Let \(1\to K\to E\xrightarrow{q}S_3\to1\) be an extension
of profinite groups with \(K\) pro-two. If a transposition
\(t\in S_3\) has an involution lift \(\widetilde t\in E\),
then there is a continuous section \(j:S_3\to E\) satisfying
\(j(t)=\widetilde t\).

\begin{proof}
For every finite \(\mathbb F_2[S_3]\)-module \(M\), restriction
\[
 H^n(S_3,M)\longrightarrow H^n(\langle t\rangle,M)
\]
is an isomorphism for \(n\geq1\). Indeed the double-coset
formula for restriction after corestriction has the identity
term and one term through the trivial intersection
\(\langle t\rangle\cap r\langle t\rangle r^{-1}=1\),
where \(r\) is a three-cycle. The second term is zero
in positive degrees. Corestriction after restriction is
multiplication by the odd index three, hence the identity
on these \(\mathbb F_2\)-vector spaces.

First suppose \(K\) is finite. Induct on its order.
When \(K\ne1\), the subgroup
\(A=\{z\in Z(K):z^2=1\}\) is a nonzero elementary
abelian characteristic subgroup of \(K\), hence normal
in \(E\). Induction gives a section of \(E/A\to S_3\)
with the prescribed image of \(\widetilde t\).
Pull back that section to an extension with kernel \(A\).
Its class in \(H^2(S_3,A)\) restricts to zero, since
\(\widetilde t\) is an involution. The proved injectivity
of restriction makes the extension split.
The difference between any splitting on \(\langle t\rangle\)
and the specified lift is an \(A\)-valued one-cocycle.
Surjectivity in degree one extends its cohomology class to
\(S_3\); changing the section by that cocycle and then by
an \(A\)-conjugation makes its value at \(t\) literally
\(\widetilde t\). This proves the finite assertion.

For general pro-two \(K\), take all open normal subgroups
\(L\) of \(E\) contained in \(K\). Each finite quotient
has such a section with the prescribed transposition.
The corresponding closed subsets of the compact space
of tuples \((e_g)_{g\in S_3}\in E^{S_3}\), imposing the
section identities modulo \(L\), are nonempty and have
the finite intersection property. Their intersection
gives a section with all identities literal.
\end{proof}

\statement{A.17}{Proved}{The parity-one relative frame can be chosen anti-invariant under conjugation}
In the notation of R2, put \(H_0=F'F^3\),
\(\Gamma_2=\widehat B/O^2(H_0)\), and let
\(d=\operatorname{Inn}(s)\circ c\), so that
\(d(s)=s\), \(d(r)=r^{-1}\), \(d(x)=y^{-1}\),
and \(d(y)=x^{-1}\). There is an actual element
\[
 b\in C_{H_0^{(2)}}(r),\qquad q(b)=[5I],\qquad d(b)=b^{-1}.
\]
Consequently the identity automorphism of
\(\widehat B/R_{\mathrm{nil}}(H_0)\), with \(s\)- and
\(\tau\)-frames one and \(r\)-frame equal to \(b\) at two
and one at all other primes, is a framed cocycle for
conjugation by \(c\). Its modulo-twelve frame parity is one.
This establishes local equivariance of these data; it does
not assert an equivariant full lift.

\begin{proof}
Use R2's actual joint quotient
\[
 J=\operatorname{im}\bigl(\widehat B\to Q\times D_0\bigr),
 \qquad Q=\mathrm{SL}_2(\mathbb Z/12)/\{\pm I\},\quad
 D_0=\widehat B/H_0.
\]
The map \(\Gamma_2\to J\) has pro-two kernel; both factors
and the joint quotient are preserved by \(d\).
The element \(z_0=([5I],1)\in J\) is central of order two
and is fixed by \(d\). Thus
\(S=\langle\bar r,z_0d\rangle\) in \(J\rtimes\langle d\rangle\)
is an actual copy of \(S_3\), with its named three-cycle
and transposition.

There is a particularly small literal lift of that
transposition as an involution. Set
\[
 C=[x,y],\qquad h=CxCx^{-1}\in F'\subseteq H_0.
\]
In R2's signed matrix coordinates modulo twelve,
\[
 q(h)=\begin{pmatrix}9&4\\8&9\end{pmatrix},\qquad
 D=\begin{pmatrix}0&1\\1&0\end{pmatrix},\qquad
 q(h)\bigl(Dq(h)D^{-1}\bigr)^{-1}=5I.
\]
Conjugation by \(D\) represents \(d\) on \(Q\).
The \(D_0\)-coordinate of \(h\) is one, so the actual
involution \(hdh^{-1}\) maps to \(z_0d\).
Apply A.16 to the preimage of \(S\) in
\(\Gamma_2\rtimes\langle d\rangle\), retaining this
involution. Its lifted three-cycle is conjugate to
the named \(r\) by the pro-two kernel, by the
coprime complement-conjugacy argument of A.9.
After that conjugation write the lifted transposition
as \(bd\). Its image in \(D_0\rtimes\langle d\rangle\)
is \(d\), so \(b\in H_0^{(2)}\), and its \(Q\)-coordinate
gives \(q(b)=[5I]\). Its order two gives \(b\,d(b)=1\).
Since it inverts the named \(r\), while \(d\) also
inverts \(r\), the element \(b\) centralizes \(r\).

All components of the prescribed relative frame have
image one in \(D_0\), so they give an actual element
of the relative fibre product. For a multiplier-one
normalized automorphism \(\Phi\) of \(\widehat B\),
write \(\Phi(s)=\beta s\beta^{-1}\) and
\(\Phi(r)=A rA^{-1}\). Direct substitution from
\(c(s)=s\), \(c(r)=sr^{-1}s\) gives the frames of
\(c\Phi c\):
\[
 \beta^c=c(\beta),\qquad
 A^c=c(\beta)s\,c(\beta)^{-1}c(A)s.
\]
For the specified local data \(\beta=1\), the second
formula is \(A^c=d(A)=A^{-1}\). These are exactly
the frames of the inverse of the identity framed
automorphism. The parity is the same as in R2's E.58:
the signed \(r\)-frame is \(5I\), whose modulo-three
value is \(-I\), and the \(s\)-frame is one.
The independently executed certificate
\texttt{check\_A17\_frame.py} verifies the literal
matrix word and the covariance identities just used.
\end{proof}

The remaining full obstruction can now be placed in a
smaller, specified kernel. Let \(L\) consist of the
multiplier-one elements whose normalized action is
identity on \(\widehat B/R_{\mathrm{nil}}(H_0)\) and whose
two pure frames both have image one there. It is the
kernel of the framed action, is stable under \(c\), and
\(\Delta(L)=0\). R2 supplies a full \(\sigma\) lifting
A.17's data, with \(\Delta(\sigma)=1\). Put \(a=\sigma c\).
Then \(a^2\in L\). A parity-one involution in this
specified fibre exists exactly when some \(v\in L\)
solves
\[
                   v\,(ava^{-1})=a^{-2}.
\]
The local construction proves that the right-hand side
lies in this framed kernel. Solvability of this last full
norm equation remains open.

\statement{A.18}{Refuted}{R2 does not admit an unrestricted order-preserving refinement}
There is an actual involution of
\(\operatorname{Out}^{\sharp}(F/R_{\mathrm{nil}}(F))\)
which lifts to a full peripheral automorphism by R2 but
has no involution lift in \(\mathcal P\).

\begin{proof}
Use the actual product
\(F/R_{\mathrm{nil}}(F)=\prod_pF^{(p)}\).
Define \(\bar a\) to be standard conjugation \(c\)
on \(F^{(3)}\) and identity on every other factor.
This is a continuous automorphism of order two,
preserving each of the three peripheral conjugacy
classes up to the common unit
\[
 \lambda_3=-1,\qquad \lambda_p=1\quad(p\ne3).
\]
Its outer class is nonidentity, since its action on
the three-primary abelianization is \(-I\).
By A.2, which applies R2 after Galois translation
of the multiplier, this class has a lift in
\(\mathcal P\). Every lift has the displayed
multiplier, because the product quotient retains
the entire abelianization \(\widehat{\mathbb Z}^2\).
By R1, a nonidentity full involution must instead
have multiplier \(-1\) at every prime. Hence no
lift of \(\bar a\) can be an involution.
\end{proof}

This counterexample is in the actual peripheral level
groups, rather than in an unrelated finite extension.
It excludes a blanket equivariant strengthening of
the central-refinement count. A.17's remaining fibre
has multiplier globally \(-1\), so this orientation
obstruction does not decide that fibre.

\statement{A.19}{Proved}{The square obstruction at an abelian refinement and the character-lifting test}
Suppose \(a\) is a lift of a quotient involution,
\(a^2\in L\), and \(L_0\triangleleft\langle L,a\rangle\)
is contained in \(L\), with \(M=L/L_0\) finite abelian.
Conjugation by \(a\) makes \(M\) a \(C_2\)-module.
An involution exists above that quotient element in
\(\langle L,a\rangle/L_0\) exactly when
\[
 [a^2]=0\quad\hbox{in}\quad
 H^2(C_2,M)=M^{a}/(1+a)M.
\]
This obstruction vanishes when \(M\) has odd order.
For elementary abelian two-primary \(M\), its target is
\(\ker(a+1)/\operatorname{im}(a+1)\).

If \(\Delta\) has a continuous lift to a character
\(\widehat{GT}_0\to\mathbb Z/4\), then it vanishes on
every involution. The same conclusion follows from
a continuous lift of \(\delta\) on \(\mathcal H\)
which is invariant under conjugation by \(c\).
Existence of either lift is not established here.

\begin{proof}
The square of the conjugation action is conjugation
by \(a^2\in L\), which is trivial on \(M\).
Also \(a^2\) is invariant under conjugation by \(a\).
Replacing \(a\) by \(va\), with \(v\in L\), changes
its square in additive notation by \((1+a)\bar v\).
This proves the quotient criterion, which is the
usual cyclic two-cocycle description of \(H^2\).
For an invariant \(u\) of odd order, \(u/2\) is
defined and \((1+a)(u/2)=u\), proving vanishing.
In characteristic two the invariants are
\(\ker(a+1)\), giving the last formula.

A character to \(\mathbb Z/4\) sends an involution
to zero or two; both reduce to zero modulo two.
Alternatively, if \(b\in\mathcal H\) is a cocycle
for \(c\) and \(\theta\) is a \(c\)-invariant
lift of \(\delta\), then
\(0=\theta(bc(b))=2\theta(b)\), so
\(\delta(b)=0\).
\end{proof}

For the actual \(a=\sigma c\) of A.17, any proposed
global character lift \(\Theta\) would necessarily satisfy
\(\Theta(a^2)=2\), although \(a^2\in L\) and
\(\Delta(L)=0\). Thus such a character must detect the
residual square inside the framed kernel; it cannot be
inferred from the local frames alone. Conversely A.17's
local parity-one framed involution rules out a character
lift to \(\mathbb Z/4\) on that local framed group itself.

\statement{A.20}{Proved}{The locally glued symmetric subgroup does not lift in the standard extension}
Let \(H_0\), \(d\), and \(b\) be as in A.17. In the standard
relative extension
\[
 \bigl(\widehat B/R_{\mathrm{nil}}(H_0)\bigr)
                  \rtimes\langle d\rangle
\]
there is a subgroup \(S_{\mathrm{loc}}\simeq S_3\), generated
by the named \(r\) and a transposition with component \(bd\)
at two and \(d\) at every odd prime. No finite subgroup
of the standard full extension
\(\widehat B\rtimes\langle d\rangle\) maps onto
\(S_{\mathrm{loc}}\). The obstruction is visible already in
\(Q\rtimes\langle d\rangle\), with \(|Q|=576\).

\begin{proof}
All local components of the transposition have the same
image \(d\) in \(D_0\rtimes\langle d\rangle\).
They therefore form an actual element of the relative fibre
product. A.17 gives its order two and its inversion of \(r\).
Their images in \(D_0\rtimes\langle d\rangle\) generate
\(S_3\), so the resulting subgroup has order exactly six.

In the full free product \(\widehat B=C_2\amalg C_3\),
every element of order three is conjugate into the named
\(C_3\), and that factor is malnormal. These are
Herfort--Ribes Theorems A' and B'.\cite{A_HR1985}
In particular \(N_{\widehat B}(\langle r\rangle)=\langle r\rangle\).
Since \(d\) inverts \(r\),
\[
 N_{\widehat B\rtimes\langle d\rangle}(\langle r\rangle)
                         =\langle r,d\rangle\simeq S_3.
\]
Every full subgroup isomorphic to \(S_3\) intersects
\(\widehat B\) in its unique subgroup of order three:
it cannot be contained in \(\widehat B\), whose finite
subgroups conjugate into its factors by the same paper's
Theorem 1. Conjugating that order-three subgroup into
\(\langle r\rangle\) proves that every full \(S_3\)
is conjugate to \(\langle r,d\rangle\).

The image of \(S_{\mathrm{loc}}\) in the actual finite
quotient \(Q\rtimes\langle d\rangle\) is
\[
                   \langle R,[5I]d\rangle .
\]
It is not conjugate there to \(\langle R,d\rangle\).
Indeed a conjugator may be taken in \(Q\), since \(d\)
normalizes the latter subgroup, and it must normalize
its unique order-three subgroup. Directly solving the
two-by-two equations in \(\mathrm{SL}_2(\mathbb F_3)\)
shows that no determinant-one matrix conjugates \(R\)
to \(R^{-1}\), and its centralizer is \(\{\pm R^j\}\).
The centralizer modulo four is also \(\{\pm R^j\}\).
Traces exclude a negative sign in the projective
commutation equation. Thus a representative of a
normalizing element has coordinates
\(\varepsilon_4R^i\) modulo four and
\(\varepsilon_3R^j\) modulo three.
Conjugation by \(D\) fixes both signs and inverts \(R\).
Its \(d\)-norm consequently has positive signs in both
coordinates. It cannot lie in \([5I]\langle R\rangle\):
one global choice of representative has negative sign
modulo three, and the other has negative sign modulo four.
This proves the nonconjugacy and hence the asserted
nonlifting of an \(S_3\) subgroup.

For completeness, a finite subgroup mapping onto
\(S_{\mathrm{loc}}\) has trivial kernel: the kernel lies
in the torsion-free free profinite group \(H_0\).
It would therefore be an \(S_3\), already excluded.
The exact exhaustive certificate
\texttt{check\_A20\_standard\_extension.py} generates all
\(576\) elements of \(Q\), obtains normalizer and
centralizer orders \(18\), lists all nine \(d\)-norm
values and all three elements of the target coset, and
checks that the lists are disjoint. It independently
enumerates the two signed centralizers used above.
\end{proof}

A.20 blocks lifting this subgroup inside the \emph{standard}
extension. A hypothetical new involution of
\(\widehat{GT}_0\) would change the action defining the
full extension. No obstruction in that unknown twisted
extension is inferred from A.20.

\statement{A.21}{Refuted}{Real projectivity and a rank-two open free subgroup do not imply finite graph recognition}
There exists a finitely generated real-projective profinite
group \(E_*\) with an open normal subgroup isomorphic to
\(\widehat F_2\), quotient \(C_2\), and self-centralizing
involutions, which is not the profinite fundamental group
of a finite graph of finite groups. It has one conjugacy
class of involutions, so it does not satisfy R1's
three-section conclusion.

\begin{proof}
Put \(A=\prod_{p\text{ odd}}\mathbb Z_p\) and define
\[
 \Gamma_*=A\rtimes\mathbb Z_2,
 \qquad u\cdot a=(-1)^{u\bmod2}a.
\]
All groups and actions here are profinite and continuous.
The kernel of the parity map \(\Gamma_*\to C_2\) is
\(A\times2\mathbb Z_2\simeq\widehat{\mathbb Z}\).
The group \(\Gamma_*\) is torsion-free: a finite-order
element projects trivially to \(\mathbb Z_2\), and then
belongs to the torsion-free group \(A\).
Set
\[
 E_*=C_2\amalg\Gamma_*,\qquad
 \varepsilon:E_*\longrightarrow C_2,
\]
where \(\varepsilon\) is identity on the first factor
and parity on the second. In the open-subgroup Kurosh
decomposition of \(\ker\varepsilon\), the two vertex
groups are \(1\) and \(\widehat{\mathbb Z}\), and the
quotient graph has two vertices joined by two edges.
Consequently
\[
 \ker\varepsilon\simeq
 \widehat{\mathbb Z}\amalg\widehat{\mathbb Z}
 \simeq\widehat F_2.
\]
The open Kurosh theorem in precisely this full-profinite
form is recalled in Herfort--Ribes, Section~2.
Their Theorems A', B', and 1 show that all nonidentity
torsion lies in conjugates of the single \(C_2\)
factor, and that each such involution has centralizer
exactly that factor.\cite{A_HR1985}
Haran's characterization of real projectivity therefore
applies: an index-at-most-two subgroup of cohomological
dimension at most one, together with self-centralizing
involutions, suffices. We use its exact formulation in
P\'al--Quick, Theorem~2.1.\cite{A_PalQuickFormal2025}

Suppose \(E_*\) were the profinite fundamental group of
a finite graph of finite groups. Its vertex groups have
order at most two. Collapse each nonloop edge whose
nontrivial edge group is equal to an endpoint group.
In the resulting reduced graph, a loop with edge group
\(C_2\) would make its stable letter centralize that
\(C_2\); the stable letter has infinite order, as is
seen by killing the vertex groups and mapping that
loop to \(\widehat{\mathbb Z}\). This contradicts the
centralizer just proved. All remaining edges therefore
have trivial group. We obtain
\[
 E_*\simeq\bigl(\mathop{\amalg}_{j=1}^{k}C_2\bigr)
                   \amalg\widehat F_r .
\]
The single torsion class forces \(k=1\). The
torsion-free index-two subgroup of this free product
has rank \(2r+k-1\), by its two-sheeted graph or by
open Kurosh. The rank two already computed forces
\(r=1\).

On the other hand \([\Gamma_*,\Gamma_*]=A\), since
multiplication by two is invertible on \(A\).
Thus
\[
 E_*^{\mathrm{ab}}=C_2\times\mathbb Z_2,
\]
whereas \((C_2\amalg\widehat{\mathbb Z})^{\mathrm{ab}}
=C_2\times\widehat{\mathbb Z}\). The latter has
nontrivial odd-primary quotients and the former does not.
This contradiction proves the claim.
\end{proof}

The construction uses the torsion-free virtually
procyclic mechanism already present in Haran's
counterexamples to full-profinite virtual freeness
recognition.\cite{A_Haran1985}
It identifies a genuine limitation of that method;
it is not a new peripheral involution. For the groups
in A.14 the three-section hypothesis, and ultimately
the marked cusp products in A.11, still have to be used.

\statement{A.22}{Proved}{The three distinguished reflections normally generate, and a proper generated subgroup has infinite index}
Let \(\alpha\) and \(h\) be as in R1, and put
\(E=F\rtimes_\alpha\langle t\rangle\), \(b=hyh^{-1}\).
The three involutions
\[
       t_0=t,\qquad t_1=xt,\qquad t_2=b^{-1}t
\]
represent the three section classes and normally
generate \(E\). If their closed generated subgroup
\(J\) is proper, then \([E:J]\) is infinite.
Normal generation alone does not supply the full
generation required in A.11.

\begin{proof}
The equality \(f=h^{-1}\alpha(h)\) gives
\(\alpha(b)=b^{-1}\), so all three displayed
elements have order two. In
\(E^{\mathrm{ab}}\simeq C_2^3\), their images are
\(t,x+t,y+t\); these are distinct and hence the
elements are pairwise nonconjugate. R1 states that
there are exactly three classes. Quotienting \(E\)
by their closed normal closure kills \(t\), then
\(x=t_1t\), then \(b^{-1}=t_2t\). Since \(b\) is a
conjugate of \(y\), it kills \(y\), and the quotient
is trivial.

There is a continuous epimorphism
\(C_2\amalg C_2\amalg C_2\to J\) sending the three
factors to the displayed involutions. Its parity
kernel is free profinite of rank two, so \(J\cap F\)
is generated by at most two elements. If \(J\)
had finite index in \(E\), then \(J\cap F\) would
be open in \(F\) and the profinite Schreier formula
would give
\[
 d(J\cap F)=1+[F:J\cap F].
\]
It follows that \(J\cap F=F\), and \(t\in J\)
then gives \(J=E\). This proves the assertion.
\end{proof}

\statement{A.23}{Proved}{Torsion and cusp centralizers in the relative primary modular groups}
Let \(q:\widehat B=C_2\amalg C_3\to D\) be a finite
quotient injective on the two factors, with torsion-free
kernel \(U\), and write
\[
 \Gamma_p=\widehat B/O^p(U),\qquad V_p=U^{(p)}.
\]
For \(a=s\) or \(r\), of prime order \(\ell\), the
following statements hold:
\[
 \begin{array}{ll}
 p\ne\ell:&C_{\Gamma_p}(a)\longrightarrow C_D(q(a))
                    \text{ is onto},\\[2pt]
 p=\ell:&C_{\Gamma_p}(a)=\langle a\rangle.
 \end{array}
\]
For \(U=U_{12}=\ker q_{12}\) in the projective modular
quotient, put \(\Gamma=\widehat B/R_{\mathrm{nil}}(U)\).
Then
\[
 C_{\Gamma_p}(\tau)=\overline{\langle\tau\rangle},
 \qquad C_\Gamma(\tau)=\overline{\langle\tau\rangle}
                    \simeq\widehat{\mathbb Z}.
\]
In the second formula the closure is taken in \(\Gamma\);
in the first it is taken in \(\Gamma_p\).

\begin{proof}
If \(p\ne\ell\), choose a lift \(g\in\Gamma_p\)
of an element of \(C_D(q(a))\). The subgroups
\(\langle a\rangle\) and
\(\langle gag^{-1}\rangle\) are complements to
\(V_p\) in the inverse image of \(\langle q(a)\rangle\).
Coprime Schur--Zassenhaus, first in finite quotients
and then by compactness, conjugates the second to
the first by an element of \(V_p\). The generators
with image \(q(a)\) then agree. Adjusting \(g\)
by that element gives the asserted centralizer lift.

Now let \(p=\ell\), and put
\(J=q^{-1}(\langle q(a)\rangle)\subset\widehat B\).
The open Kurosh decomposition expresses \(J\) as
a finite free profinite product of copies of
\(C_p\), together with a finitely generated free
profinite factor. It has no factor of the other
prime order. Since \(J/U=C_p\),
\[
          J/O^p(U)=J^{(p)}.
\]
Indeed the left side is pro-\(p\), and every
homomorphism from \(J\) to a finite \(p\)-group
kills \(O^p(U)\). Its maximal pro-\(p\) quotient
is the corresponding free pro-\(p\) product.

If \(g\in C_{\Gamma_p}(a)\), choose
\(g_0\in\widehat B\) with the same \(D\)-image.
The elements \(a\) and \(g_0ag_0^{-1}\) belong
to \(J\) and become conjugate by \(V_p\) in
\(J^{(p)}\). Distinct Kurosh factors remain
nonconjugate in that free pro-\(p\) product,
by Herfort--Ribes Theorem B'. Their double-coset
labels in \(J\backslash\widehat B/\langle a\rangle\)
must consequently agree. Thus \(g_0\in J\),
so \(g\in J^{(p)}\). The same theorem gives
\(C_{J^{(p)}}(a)=\langle a\rangle\).
This proves the prime-primary assertion without
assuming a full-profinite Kurosh theorem for
arbitrary closed subgroups.\cite{A_HR1985}

For the cusp assertion, \(|D|=576\) and
\(q(\tau)\) has order twelve. The discrete kernel
has 48 cusps, genus 25, and free rank 97:
the cusp count is \(576/12\), and its Euler
characteristic is \(576(-1/6)=-96\).
In \(H_1(U,\mathbb F_p)\), the 48 oriented cusp
vectors have precisely the relation that their
sum is zero. In particular they are pairwise
distinct for every prime, including two and three.
The \(D\)-stabilizer of the vector represented
by \(\tau^{12}\) is exactly
\(\langle q(\tau)\rangle\): deck transformations
permute these vectors by their cusp permutation.
These assertions follow directly from the usual
oriented surface presentation of this torsion-free
finite cover of the modular orbifold.

If \(g\in\Gamma_p\) centralizes \(\tau\), it
centralizes \(\tau^{12}\in V_p\). Its image in
\(D\) therefore belongs to the indicated cusp
stabilizer. After multiplying \(g\) by an integral
power of \(\tau\), it belongs to \(V_p\).
The boundary element \(\tau^{12}\) is primitive
in \(U\), since another cusp can be omitted in
the surface free basis, and remains primitive
in \(V_p\). Herfort--Ribes Theorem B' in that
free pro-\(p\) group gives
\(C_{V_p}(\tau^{12})=\overline{\langle\tau^{12}\rangle}\).
This proves the first cusp equality. The fibre
product of these procyclic groups over
\(\langle q(\tau)\rangle=C_{12}\) is exactly
\(\widehat{\mathbb Z}\): its \(p\)-coordinate
is the infinite \(p\)-primary exponent, and the
remaining finite coordinates are the common
residue modulo twelve. This proves the second.
The same proof applies to \(x=\tau^{-2}\): its
negative sixth power is \(\tau^{12}\), so a
centralizer of \(x\) fixes the same cusp vector.
Conjugating by \(s\) treats \(y=sxs\). Thus
\(C_\Gamma(x)=\overline{\langle\tau\rangle}\)
and \(C_\Gamma(y)=s\overline{\langle\tau\rangle}s\),
as used below.
\end{proof}

\statement{A.24}{Proved}{The parity character descends to the actual outer image at level twelve}
Let
\[
 F_{12}=F/R_{\mathrm{nil}}(U_{12}),\qquad
 G_{12}=\operatorname{im}\!\left(
 \widehat{GT}_0\longrightarrow
               \operatorname{Out}^{\sharp}(F_{12})\right).
\]
The map is well-defined by the finite covariance in
refereed LM7\_PUSH E.65, which preserves \(U_{12}\).
There is a unique continuous character
\(\Delta_{12}:G_{12}\to\mathbb F_2\) such that
\(\Delta=\Delta_{12}\circ\pi_{12}\).

\begin{proof}
First consider a normalized multiplier-one action
\(\Phi\) which is identity on
\(\Gamma=\widehat B/R_{\mathrm{nil}}(U_{12})\).
Its pure \(s\)-frame, when projected to \(\Gamma_2\),
centralizes \(s\), so A.23 puts it in
\(\langle s\rangle\). Purity makes this projection
identity, since \(F\cap\langle s\rangle=1\)
already in the quotient \(S_3\). Hence its
\(Q_{12}\)-image is identity. At \(p=3\) the
same argument applies to the pure \(r\)-frame.
Moreover \(\Phi(s)=sf\) gives \(q_{12}(f)=1\).
The explicit covariance formula of LM7\_PUSH
E.65.9 then gives \(\rho_\Phi\equiv0\pmod3\),
and E.65.11 gives \(\rho_\Phi\in2\widehat{\mathbb Z}\).
Thus \(V_\Phi=I\) modulo twelve. Both parity
bits are zero. This proves that \(\Delta\)
vanishes on the kernel of the actual \(B\)-action.

We check that passing to an outer action on
\(F_{12}\) adds no kernel. Its abelianization is
\(\widehat{\mathbb Z}^2\), because
\(R_{\mathrm{nil}}(U_{12})\subset U_{12}'\subset F'\).
An outer-trivial action therefore has multiplier
one. Suppose its restriction is conjugation by
\(u\in F_{12}\). The centralizer of \(F_{12}\)
in \(\Gamma\) is trivial. To see this, the cusp
argument in A.23 also gives
\(C_\Gamma(x)=\overline{\langle\tau\rangle}\)
and
\(C_\Gamma(y)=s\overline{\langle\tau\rangle}s\).
Their images in \(S_3\) are distinct order-two
subgroups; an element of their intersection lies
in \(F_{12}\), and then in both
\(\overline{\langle x\rangle}\) and
\(\overline{\langle y\rangle}\). Abelianization
shows that their intersection is trivial.

An automorphism of \(\Gamma\) fixing the open
normal subgroup \(F_{12}\) pointwise is consequently
identity: for each \(g\in\Gamma\), its discrepancy
\(\Phi(g)g^{-1}\) centralizes that subgroup.
Thus the given \(\Phi\) is conjugation by \(u\)
on all of \(\Gamma\). It fixes \(\tau\), so A.23
gives \(u\in\overline{\langle\tau\rangle}\cap F_{12}
=\overline{\langle x\rangle}\). Write \(u=x^a\).
The normalized formula \(\Phi(s)=sf\) then gives
\(f=y^a x^{-a}\). Since \(f\in F_{12}'\), its
abelian coordinates \((-a,a)\) vanish; hence
\(a=0\). The actual action was identity after all.

We have proved kernel containment for \(\Delta\).
It therefore factors through \(G_{12}\); continuity
follows because a continuous surjection of compact
groups to a Hausdorff group is a quotient map.
\end{proof}

\statement{A.25}{Proved}{An actual parity-one involution exists at the relative level twelve}
The group \(G_{12}\) of A.24 contains an involution
\(a_{12}\) with
\[
             \lambda(a_{12})=-1,
             \qquad\Delta_{12}(a_{12})=1.
\]
It is not conjugate to the image of \(c\) in
\(G_{12}\). Thus the first, specified-relative-level
assertion of Lead~3 holds at \(U_{12}\).
The existence of an involution lifting it in the
full group is not asserted.

\begin{proof}
Let \(r_0=(3/4)^{1/4}>0\). The splitting field
\(L=\mathbb Q(i,r_0)\) has degree eight: the
polynomial \(4X^4-3\) is Eisenstein at three,
and \(\mathbb Q(r_0)\) is real. Choose a Galois
element \(\gamma\) whose restriction to \(L\)
fixes \(i\) and sends \(r_0\) to \(ir_0\).
Such a restriction is an automorphism of \(L\)
and extends to \(\overline{\mathbb Q}\).
Use its actual Galois image in \(\widehat{GT}\),
and put
\[
 \sigma_0=\gamma c\gamma^{-1}c,
 \qquad \sigma=\sigma_0^3.
\]
Then \(\lambda(\sigma_0)=1\) and
\(c\sigma_0c=\sigma_0^{-1}\), so
\(\sigma c\) is an actual full involution.
Also \(\sigma_0(r_0)=-r_0\). The Galois comparison
in refereed LM7\_PUSH E.65.13, including the model
and sign calculation supplied by its referee,
makes both normalized CM bits of \(\sigma_0\) equal
to one. More precisely these are
\(\varepsilon_3(T_{\sigma_0})\) and
\(\varepsilon_2(W_{\sigma_0})\), with the shear
\(V\) included in the definitions of \(T,W\).
LM7\_PUSH E.65.12 proves that \(T,W\) are
homomorphisms on \(\mathcal H\); composing with
the two indicated characters proves that both
bits are characters on all of \(\mathcal H\).
The same bits therefore hold for \(\sigma\).

The cubing calculation in the last paragraph of
that refereed proof gives
\[
 q\Phi_\sigma=q,
 \qquad q_{12}(f_\sigma)=I,
 \qquad V_\sigma=I.
\]
Let \(A,\beta\in F\) be its actual pure frames,
and write \(a=q(A),\ b=q(\beta)\). The r-bit
and s-bit are both one, while \(a\) and \(b\)
centralize \(q(r)\) and \(q(s)\), respectively.
More precisely, cubing the two finite cyclic
characters gives
\[
 a=[5I],\qquad
 b=[B]\ \text{or}\ [B^{-1}],\qquad
 B=\begin{pmatrix}9&4\\8&9\end{pmatrix}\pmod{12}.
\]
Here \(B\equiv I\pmod4\), \(B\equiv S\pmod3\),
and \(B^2=5I\). The choice between the two
possibilities will not matter.

On \(\Gamma_p=\widehat B/O^p(U_{12})\), define
\[
 \phi_2=\Phi_\sigma,\qquad
 \phi_p=\operatorname{id}\quad(p\ne2).
\]
All induce identity on the shared finite quotient
\(Q_{12}\). They therefore define an actual
continuous automorphism \(\phi\) of their fibre
product \(\Gamma\). Each fixes \(\tau\).
The automorphism \(\phi c\) has square one:
at two it is \(\Phi_\sigma c\), and at every
other prime it is \(c\).

We now verify full peripheral gluing with literal
conjugators, the essential point of the construction.
At two, A.23 supplies \(k_2\in C_{\Gamma_2}(r)\)
with \(q(k_2)=a^{-1}\). Replace the r-frame
\(A_2\) by \(A_2k_2\). Its \(Q_{12}\)-image
is one. At every other prime take the identity
r-frame. These give a single pure r-frame of
\(\phi\), with \(Q_{12}\)-image one.

At two retain the s-frame \(\beta_2\), with
\(Q_{12}\)-image \(b\). For every odd prime,
A.23 supplies \(k_p\in C_{\Gamma_p}(s)\)
with \(q(k_p)=b\). These are frames for the
identity automorphism at those primes. Their
common image \(b\in q(F)\) makes the resulting
global s-frame pure. The \(\tau\)-frame is one
at every prime. Thus \(\phi\) preserves all
three conjugacy classes \(s,r,\tau\), with
the stated common quotient frames.

Apply R2 with these literal conjugators. It gives
an automorphism of \(\widehat B\) fixing \(\tau\)
and lifting \(\phi\), with pure frames whose
\(Q_{12}\)-images are \((1,b)\) in the order
\((r,s)\). This lift belongs to \(\mathcal H\).
Indeed write it as \(s\mapsto sf\),
\(r\mapsto f^{-1}r\). Its reduction has
\(f_2=f_\sigma\) and \(f_p=1\) at odd primes,
so its image in every primary abelianization of
\(F\) is zero. Since
\(R_{\mathrm{nil}}(U_{12})\subset F'\), the
full \(f\) lies in \(F'\). The two torsion
relations are exactly the gentle relations, and
R2 supplies invertibility.

Denote this full multiplier-one lift by
\(\widetilde\sigma\). Its action on \(Q_{12}\)
is identity and its \(f\)-image there is one;
the covariance formula makes \(V=I\). Its
r-bit is zero and its s-bit is one. Hence
\(\Delta(\widetilde\sigma)=1\). The outer
image
\[
                  a_{12}=\pi_{12}(\widetilde\sigma c)
\]
is precisely the involution already constructed.
Its action on \(F_{12}^{\mathrm{ab}}\) is \(-I\),
so it has order exactly two. A.24 gives
\(\Delta_{12}(a_{12})=1\), whereas
\(\Delta_{12}(\pi_{12}(c))=0\). The character
separates their conjugacy classes.
\end{proof}

The exact finite computations underlying the construction
are retained in \path{check_A25_level.py} and its output:
the quotient has order 576, the pure subgroup has order 96,
the cusp orbit has size 48, and the frame matrix has
square \([5I]\). The script also verifies the dihedral
action on the four roots used above. The centralizer
and lifting assertions are proved in the text, not
inferred from this finite computation.

Unlike A.17, A.25 is an actual involution in an
outer image of the relative level. It uses the
different primes where torsion centralizers acquire
extra lifts, and retains their common finite quotient
frames throughout. Let
\(L_{12}=\ker(\widehat{GT}_0\to G_{12})\) and
\(a=\widetilde\sigma c\). Then
\(a^2\in L_{12}\) and \(\Delta(L_{12})=0\).
A full involution above this specified class exists
if and only if
\[
          v(ava^{-1})=a^{-2}
          \quad\text{has a solution }v\in L_{12}.
\]
A.19 describes every abelian refinement of this
remaining square obstruction. No solution or
incompatible finite refinement has yet been proved.

\statement{A.26}{Proved}{Finite frame visibility prevents arithmetic-only continuation}
Put
\[
 U=U_{12},\qquad V=\overline{U'U^6},\qquad D_6=\widehat B/V.
\]
This is a finite quotient of order \(576\cdot6^{97}\),
and \(V\) is characteristic in \(U\), normal in
\(\widehat B\), and preserved by \(\widehat{GT}_0\).
The two normalized CM bits of an element of
\(\mathcal H\) are determined by its actual action
on \(D_6\). In particular both bits are zero for
an element acting identically on \(D_6\).

Let \(H_6\) be the actual image of \(\mathcal H\)
in \(\operatorname{Aut}(D_6)\), and
\(J_6=\langle H_6,c_6\rangle\). There is a
character \(\Delta_6:J_6\to\mathbb F_2\) induced
by \(\Delta\) on
\(\mathcal H\rtimes\langle c\rangle\).
The involution of A.25 induces an actual involution
\(a_6\in J_6\) with \(\Delta_6(a_6)=1\).
No Galois element of multiplier \(-1\) has the
same actual action on \(D_6\).

\begin{proof}
For \(p=2,3\), write
\(D_p=\widehat B/\overline{U'U^p}\).
If \(a=s\) when \(p=2\), or \(a=r\) when
\(p=3\), then
\[
  \operatorname{im}\bigl(C_{D_p}(a)\to Q_{12}\bigr)
                      =\langle q(a)\rangle.
\]
To prove this, let \(g\) be a centralizing element
and choose a lift \(g_0\in\widehat B\) of its
\(Q_{12}\)-image. In
\(J=q^{-1}(\langle q(a)\rangle)\), the elements
\(a\) and \(g_0ag_0^{-1}\) become conjugate
by an element of \(U/\overline{U'U^p}\).
If their Kurosh factor labels were different,
the free-product universal property would give
a continuous character \(J\to C_p\) sending
one to one and the other to zero. Every such
character kills \(\overline{U'U^p}\), and
characters are constant on conjugacy classes.
This is a contradiction. Therefore the labels
coincide and \(q(g_0)\in\langle q(a)\rangle\),
as claimed. The reverse inclusion is immediate.

Any two pure s-frames for the same actual action
on \(D_6\) differ by a centralizer element.
The displayed assertion at two, followed by
intersection with the pure subgroup, shows that
their \(Q_{12}\)-images agree. The assertion at
three does the same for r-frames. The shear
\(V_\Phi\) modulo twelve is determined by the
\(Q_{12}\)-image of \(f\), by the explicit
covariance formula and the evenness of \(\rho\).
Thus both normalized bits are determined.

The kernel of \(\mathcal H\to H_6\) therefore
lies in \(\ker\delta\), giving a character
\(\delta_6\) on \(H_6\). It is invariant under
\(c_6\) because \(\delta\) is c-invariant.
The element \(c_6\) is not in \(H_6\):
on the retained abelian quotient
\(\widehat B^{\mathrm{ab}}=C_6\), its action
is inversion, whereas \(H_6\) acts identically.
Hence \(J_6=H_6\rtimes\langle c_6\rangle\),
and
\(\Delta_6(hc_6^e)=\delta_6(h)\) is well-defined.

The group \(U/R_{\mathrm{nil}}(U)\) has
characteristic subgroup whose quotient is
\(U/\overline{U'U^6}\). The actual involution
\(\phi c\) in A.25 therefore induces an actual
involution on \(D_6\). It belongs to \(J_6\),
has \(\Delta_6=1\), and is nonidentity since
it inverts the abelian quotient \(C_6\).
If a Galois element of multiplier \(-1\)
induced this action, it would belong to
\(\mathcal H\rtimes\langle c\rangle\) and
its \(\Delta\)-value would be one. This
contradicts R3's Galois vanishing.
\end{proof}

This is an obstruction to continuing the
\emph{arithmetic-only construction}, not an
obstruction to lifting \(a_{12}\) at all.
At any relative level below \(V\), the common
finite action must be \(a_6\). Thus no primary
component of such a continuation can be obtained
from a full Galois element of multiplier \(-1\).
In particular one cannot choose identity-on-
\(D_6\) Galois c-cocycles with the two CM bits
still equal to one. Nongalois relative lifts,
or a solution of the norm equation in
\(L_{12}\), remain possible. This conclusion
does not identify the entire character
\(\Delta\) on the image of arbitrary multipliers
in \(\operatorname{Aut}(D_6)\); the precise
finite domain here is \(J_6\).

\statement{A.27}{Proved}{The exact finite real-reflection count and its properness condition}
Let \(\pi:E\twoheadrightarrow D\) be an epimorphism of
finite groups. Fix involutions \(d_0,d_1,d_2\in D\)
which generate \(D\), and subsets
\(T_{ij}\subseteq\pi^{-1}(d_id_j)\), for
\((i,j)=(0,1),(1,2),(2,0)\). For every subgroup
\(J\leq E\), put
\[
 I_i(J)=\{u\in J:\pi(u)=d_i,\ u^2=1\},
 \qquad
 M_{ij}(J)_{u,v}=\mathbf1_{T_{ij}}(uv).
\]
Thus \(M_{ij}(J)\) is a rectangular matrix with rows
indexed by \(I_i(J)\) and columns indexed by
\(I_j(J)\). The number of triples of involutions
in \(E\) with the specified images and pair products,
and which generate \(E\), is
\[
 \sum_{J\leq E}\mu(J,E)
       \operatorname{tr}\bigl(M_{01}(J)M_{12}(J)M_{20}(J)\bigr),
 \tag{A.27.1}
\]
where \(\mu\) is the ordinary subgroup-lattice
M\"obius function. Empty index sets contribute zero.
This count concerns actual homomorphisms from
\(C_2*C_2*C_2\), with actual quotient data.

\begin{proof}
Expanding the trace gives exactly one summand for
\((u_0,u_1,u_2)\in I_0(J)\times I_1(J)\times I_2(J)\),
and that summand is one exactly when all three
pair-product conditions hold. Their ordered product
is automatically one:
\((u_0u_1)(u_1u_2)(u_2u_0)=1\).
Every counted triple has a unique generated subgroup.
M\"obius inversion on the finite subgroup lattice
therefore gives (A.27.1). The assumption that the
\(d_i\) generate makes every subgroup not mapping
onto \(D\) contribute zero.
\end{proof}

The matrices in (A.27.1) retain the common reflection
at the ends of adjacent edges. The ordinary
boundary-class trace sums of R2 do not retain this
incidence datum. Equality of those ordinary sums
therefore does not by itself establish equality of
(A.27.1). A concrete perfect-kernel control is
\(S_5\twoheadrightarrow C_2\), with kernel \(A_5\),
and all three \(d_i\) nonidentity. Every \(d_i\)
has an involution lift, and \(S_5\) has three
ordinary generators, but three such lifts cannot
generate \(S_5\): they are transpositions, whose
three edges cannot connect a graph on five vertices.
The certificate \path{check_A27_real_count.py}
checks perfectness of the kernel and enumerates
all 1000 triples. The generated-subgroup orders
and multiplicities are
\[
 (2,10),\quad(4,90),\quad(6,240),\quad
 (12,180),\quad(24,480).
\]
This is a control on a proposed general real
embedding theorem, not a counterexample to the
negative pants problem. Indeed a product of two
odd involutions in \(S_5\) has order at most three;
no generating pants triple of \(A_5\) has all
three orders at most three. The same certificate
checks all 816 candidate ordered pairs and finds
no pair generating \(A_5\).

\statement{A.28}{Proved}{A negative finite peripheral action is stopped by a central spin obstruction}
There is an epimorphism \(q:F\twoheadrightarrow A_6\)
and an involutory automorphism \(\vartheta\) of
\(A_6\) which sends each of the three labelled
cusp images to an \(A_6\)-conjugate of its inverse,
but no multiplier-minus-one peripheral automorphism
of \(F\) satisfies
\[
                       q\Phi=\vartheta q.
\]
In particular no full negative peripheral involution
induces this finite action. The obstruction is
already visible in a central double cover; it is
not a counterexample to R2 or A.3.

\begin{proof}
We use rightmost-first permutation composition and put
\[
\begin{aligned}
 X_0&=(1\,2\,3\,4)(5\,6),&
 Y_0&=(1\,2\,4\,6)(3\,5),\\
 Z_0&=(1\,2\,6\,3)(4\,5),&
 t&=(2\,4).
\end{aligned}
\]
The certificate \path{check_A28_s6_passport.py}
verifies \(X_0Y_0Z_0=1\) and
\(\langle X_0,Y_0\rangle=A_6\), giving
\(q(x)=X_0,q(y)=Y_0,q(z)=Z_0\).
Let \(\vartheta=\operatorname{Ad}(t)\).
The required inverse-class conjugators are,
in the displayed cusp order,
\[
                    1,\qquad(2\,6\,4),\qquad(2\,3\,4).
\]
In particular these conjugators are even.
Each cusp also has an odd involution which
inverts it: respectively \((2\,4),(2\,6),(2\,3)\).
The permutation \((3\,4\,6)\), of order three
in \(A_6\), cyclically permutes the three
cusp images by conjugation.

For completeness there is a short section-incidence
explanation of why an order-two full lift is impossible.
Odd involutions of \(S_6\) have just two cycle types,
\((2)\) and \((2,2,2)\). The product of two of
one type cannot have cycle type \((4,2)\).
For transpositions its order is at most three.
For two fixed-point-free involutions, their two
perfect matchings decompose into alternating cycles;
on six letters the possible product orders are
again at most three. Thus every \((4,2)\) cusp
requires opposite types at its two reflection ends.

In a hypothetical full negative pants extension,
R1 gives exactly three section classes. Their
images in
\(E^{\mathrm{ab}}=C_2^3\) are
\(\bar t,\bar x+\bar t,\bar y+\bar t\).
The x-cusp joins the first two, the y-cusp the
first and third, and the z-cusp the last two:
its two ends differ by \(\bar x+\bar y\).
Here each cusp has genuine involution ends.
Indeed an element of the odd coset which inverts
a full cusp has square in its procyclic centralizer;
inversion fixes that square, so the square is one.
Their three finite class labels would therefore
colour a triangle with two colours, with opposite
colours on each edge. This is impossible. No
assumption that the three reflections generate
is made in this argument.

A stronger obstruction excludes even infinite-order
full peripheral lifts of multiplier \(-1\).
We give a finite certificate rather than assume
that an arbitrary finite peripheral automorphism
lifts. The two degree-80 permutations in the
ATLAS representation \texttt{2S6G1-p80B0}
are opened source data. The code
\path{check_A28_double_cover.py} enumerates their
joint group with the natural standard generators
\((1\,2),(2\,3\,4\,5\,6)\) of \(S_6\).
It verifies an epimorphism
\[
 \widetilde S\twoheadrightarrow S_6,
 \qquad |\widetilde S|=1440,
 \qquad \ker=\langle\epsilon\rangle\cong C_2,
\]
with central kernel. Its even preimage
\(\widetilde A\) has order 720. Choose arbitrary
lifts \(X,Y\in\widetilde A\) of \(X_0,Y_0\),
and put \(Z=(XY)^{-1}\). Choose an involution
\(T\in\widetilde S\) above \(t\).
The enumeration of all conjugators in
\(\widetilde A\) proves, for each
\(Q\in\{X,Y,Z\}\),
\[
 TQT^{-1}\not\sim_{\widetilde A}Q^{-1},
 \qquad
 TQT^{-1}\sim_{\widetilde A}\epsilon Q^{-1}.
 \tag{A.28.1}
\]
Each lift has order eight. More precisely the
unit exponents \(k\pmod8\) for which
\(TQT^{-1}\sim Q^k\) are exactly \(3,5\).
The same statements hold for either central
choice of each lift. All 1440 elements, the
central kernel, the 720 even conjugators and
the stated class tests are checked by the code;
the group order is not inferred solely from
the ATLAS label.\cite{AAtlasA6}

Let \(q_2:F\to\widetilde A\) send \(x,y\)
to \(X,Y\). If \(\Phi\) as in the statement
existed, the two homomorphisms
\(q_2\Phi\) and \(\operatorname{Ad}(T)q_2\)
would have the same projection to \(A_6\).
Their ratio is consequently a continuous central
character \(\eta:F\to\langle\epsilon\rangle\).
The multiplier-minus-one cusp conditions and
(A.28.1) require
\[
             \eta(x)=\eta(y)=\eta(z)=\epsilon.
\]
But \(xyz=1\), contradicting
\(\eta(x)\eta(y)\eta(z)=1\).
This proves the stronger assertion.
\end{proof}

The kernel of \(q\) is open normal; it is not
asserted to be characteristic. Any characteristic
refinement retaining this quotient and this action
inherits the displayed obstruction. Conversely,
we have not constructed an automorphism of such a
characteristic refinement or a relative pronilpotent
level inducing \(\vartheta\). The central calculation
explains why the finite inverse-class data alone
cannot be inserted into R2. The dyadic class
obstruction, rather than perfectness alone, is
what is retained here.

\statement{A.29}{Proved}{The new relative involution has exactly three section classes}
For the involution \(a_{12}\) constructed in A.25,
choose its actual representative \(\phi c\) on
\(F_{12}=F/R_{\mathrm{nil}}(U_{12})\).
Then
\[
                \bigl|H^1(C_2,F_{12})\bigr|=3,
\]
where the action is this specified representative.
Its relative dyadic component likewise has exactly
three section classes and has trivial fixed subgroup
in its orientation kernel. These are section
classes of the group extension, not conjugacy
classes of peripheral automorphisms.

\begin{proof}
First let
\(W=C_2\amalg C_2\amalg C_2=F\rtimes\langle c\rangle\),
and let \(U\) be any open normal subgroup of
\(W\) contained in \(F\). In
\(W/O^2(U)\), every nonidentity torsion element
of order two has nonidentity image in \(W/U\),
because \(U^{(2)}\) is torsion-free. If that
image generates \(C\cong C_2\), let
\(J\leq W\) be its inverse image. The open
Kurosh theorem expresses \(J\) as a free
profinite product of order-two factors and a
free profinite group. Since \(J/U=C_2\),
\[
                     J/O^2(U)=J^{(2)}.
\]
Indeed the left side is pro-2, and every finite
2-group quotient of \(J\) kills \(O^2(U)\).
The torsion theorem for the resulting free pro-2
product shows that the given involution is conjugate
to an image of a full factor of \(W\).
The three full factors remain nonconjugate,
because \(W^{\mathrm{ab}}=C_2^3\) survives this
quotient. Conjugacy under the even subgroup
is the same as full conjugacy for odd involutions:
one may multiply an odd conjugator by its
source involution to make it even.

The same open-product argument proves that every
such torsion centralizer is exactly the cyclic
subgroup it generates. Explicitly, if a centralizing
element has finite image \(g\), conjugating its
factor by a lift of \(g\) gives two factors of
\(J\) that become conjugate in \(J^{(2)}\).
The free pro-2 torsion theorem forces equality
of their Kurosh labels, so \(g\in C\).
The centralizer computation now takes place inside
\(J^{(2)}\), where the free-factor centralizer
is that factor. Its intersection with the even
subgroup is trivial. This uses the free-factor
case of Herfort--Ribes, not a closed-subgroup
Kurosh assertion.\cite{A_HR1985}

Apply this to \(U=U_{12}\), which is preserved
by the actual Galois action used in A.25.
The dyadic involution \(\Phi_\sigma c\) is
Galois-conjugate to \(c\), so its section set
has the same three classes and trivial even
centralizers. The choice of conjugating Galois
element preserves \(U_{12}\), by R3's quotient
covariance.

At each odd prime the kernel above the common
finite quotient is pro-\(p\). Every involution
in the specified quotient coset has an involution
lift, and any two such lifts with the same finite
image are conjugate by that kernel, by
Schur--Zassenhaus. Moreover the centralizer of
an involution surjects onto the centralizer of
its finite image: conjugate a lift by any chosen
finite centralizer lift, then apply the same
complement-conjugacy assertion in the kernel.

It follows that projection to the dyadic component
induces a bijection on section-conjugacy classes
for the complete fibre product. Surjectivity
follows by choosing odd-prime lifts with the
same finite image. For injectivity, choose a
dyadic conjugator and its finite image; at every
odd prime use centralizer surjectivity to adjust
a conjugator to this very same finite image.
The resulting tuple is an element of the fibre
product and conjugates the two sections. This
proves the asserted cardinality three.
\end{proof}

The unmarked pro-2 recognition theorem is fully
consistent with A.29. Herfort--Zalesskii's
Theorem~6.2 conjugates a finite 2-subgroup of
\(\operatorname{Aut}(F^{(2)})\) into the
automorphisms of a dense abstract free group.
It does not retain three prescribed peripheral
classes.\cite{AHerfortZalesskii2013}
Thus neither that theorem nor the section count
proves the peripheral conjugacy statement left
open in A.12. The actual involution of A.25
already passes the three-section test.

\statement{A.30}{Proved}{The level-twelve frames can also satisfy the involution relations}
The construction in A.25 can be made with its literal
relative frames \(A,\beta\in\Gamma\) so that the
associated r-normalizer and s-normalizer in
\(\Gamma\rtimes\langle\phi c\rangle\) are actual
subgroups \(S_3\) and \(C_2\times C_2\), with the
finite frame values still
\[
                         q(A)=1,\qquad q(\beta)=b.
\]
Consequently an R2 lift \(a=\widetilde\sigma c\)
can be chosen with
\[
 a^2\in L_{12}^{\mathrm{fr}},\qquad
 L_{12}^{\mathrm{fr}}
 =\{k\in\mathcal H:\Phi_k|_\Gamma=1,
                  \ A_k|_\Gamma=\beta_k|_\Gamma=1\}.
 \tag{A.30.1}
\]
This is a smaller kernel than the unframed
\(L_{12}\). It does not assert a lift at a
smaller characteristic level or solve the remaining
norm equation.

\begin{proof}
Write \(a_2=\Phi_\sigma c\) for the genuine
Galois involution used at the dyadic component,
and write \(A_2,\beta_2\) for its underlying
multiplier-one element's frames. Put
\[
 J_2=\Phi_\sigma(s)A_2,
 \qquad D_2=J_2^{-1}a_2,
 \qquad T_2=\beta_2^{-1}a_2.
\]
The element \(D_2\) inverts \(r\), because
\(c(r)=sr^{-1}s\). In the full Galois extension,
its square lies in \(C_{\widehat B}(r)=\langle r\rangle\)
and is fixed by inversion; hence \(D_2^2=1\).
Similarly \(T_2\) centralizes \(s\) and its
square lies in \(\langle s\rangle\). That square
cannot be \(s\): the abelianization of
\(\widehat B\rtimes\langle\Phi_\sigma c\rangle\)
is \(C_2^2\), and \(s\) has nonzero image.
Thus \(T_2^2=1\). These equations descend to
\(\Gamma_2\rtimes\langle a_2\rangle\).

In the common finite extended quotient, put
\(d=sc\) and \(z_0=[5I]\). Since
\(q\Phi_\sigma=q\), the frame values give
\[
                       q(D_2)=z_0d,
                       \qquad q(T_2)=b^{-1}c.
\]
Choose a pure \(h\in\Gamma_2\) with finite
image \([B]\), where \(B\) is the matrix of
A.25. Direct multiplication gives
\[
             d(B)=B^{-1},\qquad B^2=5I.
\]
It follows that the involution \(hD_2h^{-1}\)
has finite image \(d\). The inverse image of
\(\langle q(r),d\rangle\cong S_3\) in this
relative extension has pro-2 kernel. A.16 supplies
a section retaining this exact involution. Its
order-three subgroup is conjugate to \(\langle r\rangle\)
by an element of \(U_{12}^{(2)}\), by coprime
complement conjugacy. After this conjugation the
section has an involution \(D'_2\) satisfying
\[
        D'_2r(D'_2)^{-1}=r^{-1},\qquad q(D'_2)=d.
\]
Then \(k_2=D_2D'_2\) centralizes \(r\), belongs
to \(\Gamma_2\), and has finite image \(z_0\).
Replace \(A_2\) by \(A'_2=A_2k_2\).
The r-frame now has finite image one, and its
associated r-normalizing reflection is \(D'_2\).
At every odd prime retain the identity r-frame
and the standard r-normalizing reflection \(sc\).
Thus the r-frames glue, and their normalizers
glue to an actual finite \(S_3\).

At two retain \(\beta_2\) and \(T_2\).
For an odd prime \(p\), A.23 gives a surjection
\[
             C_{\Gamma_p}(s)\longrightarrow C_{Q_{12}}(q(s))
\]
with pro-\(p\) kernel, equivariant for \(c\).
The element \(b^{-1}c\) has order two in the
finite semidirect product, since \(c(b)=b^{-1}\).
Schur--Zassenhaus supplies an involution above it
in \(C_{\Gamma_p}(s)\rtimes\langle c\rangle\).
Write this involution as \(\beta_p^{-1}c\).
Then \(q(\beta_p)=b\), and \(\beta_p\) is a
frame for the identity action at that prime.
These frames glue with \(\beta_2\); their
involutions glue to \(T=\beta^{-1}(\phi c)\),
which centralizes \(s\). Since \(T\) is outside
\(\Gamma\), it and \(s\) generate \(C_2^2\).
The unchanged \(\tau\)-frame is one.

Apply R2 to these actual frames, and call the full
lift \(a=\widetilde\sigma c\). It still has
\(\Delta(a)=1\) and the same relative action as
A.25. For its s-frame, applying \(a\) twice
gives the pure conjugator \(a(\beta)\beta\).
Its relative image is one, exactly because
\((\beta^{-1}(\phi c))^2=1\).
For the r-frame put \(J=\Phi_{\widetilde\sigma}(s)A\).
Applying \(a\) twice gives the pure conjugator
\(a(J)J\), whose relative image is one because
\((J^{-1}(\phi c))^2=1\). Uniqueness of the full
pure frames identifies these with the frames of
\(a^2\). Together with its identity relative
action, this proves (A.30.1).
\end{proof}

The kernel \(L_{12}^{\mathrm{fr}}\) is invariant
under conjugation by this chosen \(a\). One may
verify this by composing actual peripheral actions
and their literal conjugators: an identity action
with both conjugators equal to one remains such
under framed conjugation. In the specified
\emph{framed} fibre, an order-two full lift exists
exactly when
\[
             v(ava^{-1})=a^{-2}
             \quad\text{for some }v\in L_{12}^{\mathrm{fr}}.
 \tag{A.30.2}
\]
Solving (A.30.2) would solve Lead~3. Failure to
solve it would exclude only this framed fibre;
a full involution over the unframed \(a_{12}\)
could still use different relative frames, and
is governed by A.25's equation in \(L_{12}\).
Neither equation is declared solvable or insoluble.

\Part{B}{The icosahedral centralizer}
\input{sections/part_B.tex}
\Part{C}{The punctured torus}
% Original successor text.  The D_*.tex files in this directory are audit extracts,
% not successor proofs.  References use the accompanying CE_references.bib.
\section{The punctured torus: what the fixed-group reduction does and does not imply}

We use continuous homomorphisms throughout.  A free profinite group means
the full completion, and a free pro-$p$ group is explicitly so designated.
For a profinite group $H$, write $H^{(p)}$ for its maximal pro-$p$ quotient
and $O^p(H)$ for the kernel of this quotient.  Cohomology of a profinite
group with finite coefficients means continuous cohomology.

\statement{C.1}{Proved}{A nontrivial pro-odd kernel in a projective group cannot be small}
Let $B$ be a closed subgroup of a free profinite group.  Suppose that
\[
 1\longrightarrow K\longrightarrow B
       \xrightarrow{\chi}\widehat{\mathbb Z}\longrightarrow1
\]
is exact.  For every prime $p$, the group $K^{(p)}$ is either trivial
or an infinitely generated free pro-$p$ group.  Consequently, if $K$ is
pro-odd and nontrivial, it has an infinitely generated free pro-$p$
quotient for some odd prime $p$.  In particular such a $K$ is not
topologically finitely generated.  If $K$ is abelian, then $K=1$,
without any pro-odd assumption.

\begin{proof}
A closed subgroup of a free profinite group has cohomological dimension
at most one at every prime and is projective in the category of profinite
groups.  These implications, including the Sylow criterion for
cohomological dimension, are recalled in Section~1 of
\cite{CEEnglerHaranKochloukovaZalesskii2004}.  The maximal pro-$p$
quotient of a projective profinite group is projective in the pro-$p$
category: lift a finite embedding problem with pro-$p$ target through the
original group; its lift factors through the maximal pro-$p$ quotient.
It is therefore free pro-$p$.  This applies to $K$ and to every closed
subgroup of $B$.

Regard $\mathbb Z_p$ as the $p$-coordinate subgroup of
$\widehat{\mathbb Z}$, with zero at all other primes, and put
$S=\chi^{-1}(\mathbb Z_p)$.  We claim that
\[
 O^p(S)=O^p(K),\qquad
 1\longrightarrow K^{(p)}\longrightarrow S^{(p)}
       \longrightarrow\mathbb Z_p\longrightarrow1.             \tag{C.1.1}
\]
The characteristic subgroup $O^p(K)$ is normal in $S$.  Its quotient
in $S$ is an extension of the pro-$p$ group $K^{(p)}$ by
$\mathbb Z_p$, hence is pro-$p$.  Thus $O^p(S)\subseteq O^p(K)$.
Conversely, restriction to $K$ of every finite $p$-group quotient of
$S$ kills $O^p(K)$, proving the reverse inclusion and the exact
sequence.  This argument proves the needed injectivity on the left;
it does not assume that maximal pro-$p$ quotients preserve arbitrary
short exact sequences.

Both nonzero groups in (C.1.1) are free pro-$p$.  If $K^{(p)}$ were
nontrivial and finitely generated, Theorem~2 of
\cite{CEEnglerHaranKochloukovaZalesskii2004}, applied with
$G=S^{(p)}$ and $N=K^{(p)}$, would give
\[
 \operatorname{vcd}_p(G/N)=\operatorname{cd}_p(G)-1=0.
\]
But $G/N=\mathbb Z_p$ has virtual cohomological dimension one: every
open subgroup is again $\mathbb Z_p$.  This contradiction proves the
first assertion.

If $K$ is pro-odd and nontrivial, it has a nontrivial finite quotient
$D$ of odd order.  The odd-order theorem
\cite{CEFeitThompson1963} says that $D$ is solvable.  A nontrivial
finite solvable group has nontrivial abelianization, hence has a
quotient $C_p$ for an odd prime $p$.  The composite $K\to C_p$
shows that $K^{(p)}\ne1$.  A quotient of a finitely generated profinite
group is finitely generated, proving the stated consequence.

Finally, if $K$ is abelian, each free pro-$p$ quotient $K^{(p)}$ is
abelian and therefore has rank at most one.  The first assertion
forces all these quotients to be trivial.  A nontrivial profinite
abelian group has a nontrivial cyclic prime-order quotient, so $K=1$.
\end{proof}

\statement{C.2}{Refuted}{The coarse structural data of LM4 D.78 do not force its kernel to vanish}
There is a two-generated projective profinite group $B$, embedded as a
retract of $\widehat F_2$, with a split epimorphism
$\chi:B\to\widehat{\mathbb Z}$ such that
\[
 B^{\mathrm{ab}}\cong\widehat{\mathbb Z},\qquad
 \dim_{\mathbb F_2}H^1(B_0,\mathbb F_2)=1
       \quad\text{for every open }B_0\leq B,
\]
but $\ker\chi$ is a nontrivial pro-$3$ group.  Thus projectivity,
these mod-two ranks, and the split pro-odd-kernel description alone
do not imply $\ker\chi=1$.  This is not an example of a torus
involution or a refutation of the actual question C.1 of the prompt.

\begin{proof}
Let $L$ be free pro-$3$ on $x,y$, let
$Z_{(3')}=\prod_{p\ne3}\mathbb Z_p$, and let $z$ be its
topological generator with every coordinate equal to one.  The map
$z\mapsto\sigma$, where $\sigma(x)=x$ and $\sigma(y)=y^{-1}$,
defines a continuous action through the quotient
$Z_{(3')}\to C_2$.  Set
\[
 B=L\rtimes Z_{(3')},\qquad a=xz.
\]
The factors $x$ and $z$ commute.  If $e_3\in\widehat{\mathbb Z}$
is the primary idempotent at three, then
$a^{e_3}=x$ and $a^{1-e_3}=z$.  Consequently $a,y$ generate $B$,
and $A=\overline{\langle a\rangle}$ is isomorphic to
$\widehat{\mathbb Z}$.

The homomorphism $L\to\mathbb Z_3$ with $x\mapsto1,y\mapsto0$
is invariant under $Z_{(3')}$.  Together with the quotient to
$Z_{(3')}$ it defines the continuous split epimorphism
\[
 \chi:B\longrightarrow\mathbb Z_3\times Z_{(3')}
                  =\widehat{\mathbb Z},\qquad \chi(a)=1.
\]
Its kernel is $\ker(L\to\mathbb Z_3)$ and contains $y\ne1$.
In the abelianization the relation $zyz^{-1}=y^{-1}$ kills $y$,
because multiplication by two is invertible in a pro-$3$ abelian
group.  There is no relation on $x$ or $z$.  This proves
$B^{\mathrm{ab}}\cong\mathbb Z_3\times Z_{(3')}$.

For an open subgroup $B_0$ the restriction of the quotient to
$Z_{(3')}$ has pro-$3$ kernel $B_0\cap L$ and an open image
$Z_0\leq Z_{(3')}$.  Every homomorphism $B_0\to C_2$ kills
this kernel.  The open procyclic group $Z_0$ has its two-primary
factor isomorphic to $\mathbb Z_2$.  Hence
\[
 \operatorname{Hom}_{\mathrm{cont}}(B_0,C_2)
 =\operatorname{Hom}_{\mathrm{cont}}(Z_0,C_2)\cong C_2.
\]

A Sylow pro-$3$ subgroup of $B$ is $L$.  For $p\ne3$, a Sylow
pro-$p$ subgroup maps isomorphically to the Sylow subgroup
$\mathbb Z_p$ of the quotient: the kernel is pro-$3$, and
surjectivity on Sylow subgroups follows from the profinite Sylow
theorem.  Thus every Sylow subgroup has cohomological dimension
one at its prime.  The criterion cited in C.1 gives
$\operatorname{cd}B\leq1$, so $B$ is projective.  The epimorphism
$\widehat F_2\to B$ defined by the two generators $a,y$ therefore
has a continuous section.  This section realizes $B$ as a closed
retract of $\widehat F_2$.

The script \path{code/check_kernel_model.py} checks the finite
quotients $(C_{3^n})^2\rtimes C_{2^m}$ for $(n,m)=(1,1),(1,2),(2,2)$.
It checks the group law, the epimorphism $\chi$, the two generators,
the derived group, and the invariant-vector witness.  These finite
checks confirm the stated finite models; projectivity and the
assertion for every open subgroup are proved above, independently
of any finite computation.
\end{proof}

\statement{C.3}{Proved}{The exact finite-coefficient witness for the actual torus kernel}
Let $E=F\rtimes\langle v\rangle$ be the full negative torus
extension of LM4 D.78, and let $T=\langle w\rangle$ be its
noncusp section class.  Suppose that the actual fixed subgroup
$B=C_F(w)$ has the split epimorphism
$\chi:B\to\widehat{\mathbb Z}$ of that row, with specified
section $t\mapsto a^t$, and set $A=\overline{\langle a\rangle}$,
$K=\ker\chi$.  For a finite continuous left $E$-set $X$, put
$Q_X=\mathbb F_2[X^T]$, with the induced $B$-action.  Then
\[
 K=1\quad\Longleftrightarrow\quad
 Q_X^B\longrightarrow Q_X^A\text{ is onto for every such }X.
                                                                  \tag{C.3.1}
\]
If $1\ne k\in K$, one can take $X=E/(ATN)$, where $N$ is
an open normal subgroup of $E$ contained in $F$ and chosen so
that $k\notin ATN$.  The identity-coset basis vector is an explicit
witness to failure of surjectivity.

\begin{proof}
The implication from left to right is immediate from $B=A$.
For the converse, $A$ centralizes $T$, and $A\cap K=1$ by
the specified splitting.  Thus $AT$ is a closed subgroup and
$k\notin AT$.  Compactness separates the point $k$ from $AT$
in some finite quotient of $E$, giving the required $N$.
The subgroup $ATN$ is open.  Its identity coset is fixed by
$A$ and by $T$, so its basis vector belongs to $Q_X^A$.
The element $k$ moves that vector.  It consequently does not
belong to $Q_X^B$, and the map in (C.3.1), being the inclusion
of fixed subspaces, is not onto.
\end{proof}

\statement{C.4}{Open}{The actual kernel in prompt C.1}
In the actual torus setting of C.3, prove or refute $K=1$.
The audited structural conclusions of LM4 D.78 imply that $K$ is
pro-odd.  C.1 shows that any counterexample must have an infinitely
generated free pro-$p$ quotient for an odd prime.  It would suffice
to prove that the actual $K$ is finitely generated, or that it is
abelian.  Neither assertion is established here.  C.2 shows exactly
why the abstract projectivity and mod-two rank arguments do not
establish either assertion.  The extra peripheral duality of the
torus extension has not been realized in the countermodel of C.2.


\InputIfFileExists{work_CE/C_audit.tex}{}{}
\InputIfFileExists{work_CE/C_quartic.tex}{}{}
\InputIfFileExists{work_CE/C_further.tex}{}{}
\Part{D}{The level-zero comparison}
\InputIfFileExists{work_DE/section_D.tex}{}{}
\Part{E}{Real structures, the five-point census, and corrections}
\section{Genus-three real structures and the congruence-kernel fibre}

Write $\Gamma_g$ for the orientation-preserving mapping class group of
a closed oriented surface $S_g$, $\Gamma_g^{\pm}$ for its extension
by orientation reversal, $P_g=\widehat\Gamma_g$ for the full profinite
completion, and $Q_g=\check\Gamma_g$ for the closure of its action
in $\operatorname{Out}(\widehat{\pi_1(S_g)})$.  Superscripts $\pm$
on completions have the same orientation meaning.  The congruence
kernel is
\[
 K_c=\ker(q:P_3\longrightarrow Q_3).
\]
In particular $K_c$ is not the kernel of a homology representation.
We use $a=0$ for a dividing real structure and $a=1$ for a
nondividing one; $r$ is its number of ovals.

\statement{E.1}{Proved}{The compact real-type injection of LM6 C.13}
For $g\geq2$, the topological types of orientation-reversing
involutions of $S_g$ inject into conjugacy classes in
$\operatorname{Out}(\widehat{\pi_1(S_g)})$.  They consequently
inject into conjugacy classes in $Q_g^{\pm}$ and $P_g^{\pm}$.
For genus three the six types are
\[
 (r,a)=(0,1),(1,1),(2,1),(3,1),(2,0),(4,0).                 \tag{E.1.1}
\]

\begin{proof}
Put $\Pi=\widehat{\pi_1(S_g)}$.  The group $\Pi$ is centre-free.
An outer action of $C_2$ on $\Pi$ therefore determines its extension
uniquely: it is the pullback of
$\operatorname{Aut}(\Pi)\to\operatorname{Out}(\Pi)$, with its
kernel identified with $\Pi$ by inner automorphisms.  For a
geometric involution $\tau$ this is the completed orbifold group
of $S_g/\tau$, with its orientation subgroup $\Pi$.

We justify the finite-order comparison used for this particular
orbifold group.  Surface groups are good: after passing to a
finite cover to trivialize any given finite coefficient module,
degree-one classes are killed by the finite-index kernels of
their characters, and degree-two classes are killed by a further
cover of degree divisible by the exponent of the coefficient
group.  These comparisons, or the equivalent cohomological
effacement criterion, give the completion comparison in degrees
zero, one and two; higher degrees vanish.  The same argument
works for the orientation double of a nonorientable surface.
A finite extension preserves goodness, by the two
Hochschild--Serre spectral sequences and exactness of completion
for a finite-index subgroup.  The groups here have finite virtual
cohomological dimension two and finite cohomology in each degree
with finite coefficients.  Filtered unions of finite coefficient
modules give the comparison for all discrete continuous
$\mathbb F_2$-modules.  Thus all hypotheses of Theorem~A of
\cite{CEBoggiZalesskii2025} apply, at $p=2$, to this single
involution and its orbifold group.

The discrete order-two subgroups correspond exactly to the
ovals.  One way to see this is to choose an invariant hyperbolic
metric: a lift of a reflection has a geodesic axis, whose image
is an oval, and all lifts over the same oval are conjugate by
the surface group.  A lift over an oval has centralizer
$\mathbb Z\times C_2$, with the infinite cyclic factor generated
by a circuit of that oval.  For a fixed-point-free involution
the quotient is a closed nonorientable surface and its group
is torsion-free.  Theorem~A therefore identifies the completed
section classes with the $r$ ovals, and identifies their
centralizers in $\Pi$ with the full closures
$\widehat{\mathbb Z}$ of these cyclic groups.  Conjugacy by the
orbifold group gives the same section classes as conjugacy by
its orientation subgroup, since multiplying a conjugator by
the source reflection changes its orientation without changing
the conjugacy.

The image $w_i\in H_1(\Pi,\mathbb F_2)$ of a generator of the
$i$th oval centralizer is well defined: any other topological
generator differs by a profinite unit, whose reduction modulo
two is one.  The ovals divide the surface exactly when
\[
                         \sum_{i=1}^r w_i=0.                 \tag{E.1.2}
\]
Indeed a mod-two two-chain with that boundary has constant
coefficient on each component of the complement.  Its boundary
is all the ovals exactly when the two sides can be consistently
coloured with opposite coefficients.  This is precisely the
separating condition.  Equivalently, in the nondividing case
the connected complement would force one coefficient on both
sides of every oval.  If $r=0$, the structure is declared
nondividing, and this case is already detected by the absence
of sections.

Thus an isomorphism of the completed outer extensions preserves
$r$ and $a$.  Conversely those invariants classify the
topological involution: the quotient is a compact surface with
$r$ boundary components, orientable of genus
$(g+1-r)/2$ in the dividing case and nonorientable of genus
$g+1-r$ in the nondividing case.  Its orientation double in the
interior, completed along the mirror boundary, recovers the
surface with involution.  The classification of compact surfaces
therefore gives uniqueness.  In genus three these formulas give
exactly (E.1.1).  Conjugacy in either mapping class completion
would imply conjugacy of the outer actions, proving the two
remaining injections.
\end{proof}

\statement{E.2}{Proved}{Marked incidence and the simultaneous orientation in LM6 C.14}
The injection of E.1 holds for individually labelled real markings,
and for unordered markings allowing nonreal conjugate pairs.
A peripheral isomorphism of the completed real surface extensions
recovers the cyclic incidence of real marked points on each oval.
For dividing surfaces, the cyclic orders on labelled marked ovals
are identified only by their simultaneous reversal.

\begin{proof}
Let $U=S_g\setminus D$.  Filling the punctures is the quotient
of the completed geometric fundamental group by the closed
normal subgroup generated by its peripheral subgroups.  An
isomorphism preserving these conjugacy classes therefore induces
the compact comparison of E.1.  At a real puncture the local
extension is $\widehat{\mathbb Z}\rtimes_{-1}C_2$.  Its two
local section classes are distinguished by the parity of the
exponent in $c^t s$.  Their maps to global section classes say
which real intervals end at this puncture.  Consequently the
local-to-global incidence graph has real intervals as vertices,
real punctures as edges, and these two section classes as flags.
Unpunctured ovals give isolated compact components.  A punctured
oval gives its cyclic incidence graph.  Retaining flags is
essential even for a one-edge loop or a two-edge cycle.  This
description proves the claimed separate unoriented cyclic
orders directly from the extension.

We prove that the signs on marked ovals have their asserted
integral meaning; a procyclic generator cannot be replaced by
an arbitrary profinite unit.  Fix a marked oval, a base interval
next to a mark, and the corresponding pointed reflection $s$
after filling.  By E.1 its geometric centralizer is
$H=\overline{\langle a\rangle}\cong\widehat{\mathbb Z}$,
where $a$ makes one oriented circuit of the actual oval.
For an open normal $s$-invariant subgroup $N\leq\Pi$, the
subgroup $N\langle s\rangle$ defines the finite real cover with
the selected real base point.  Path lifting identifies the
fibre on the selected oval with
\[
 \langle a\rangle N/N=HN/N=H/(H\cap N).                    \tag{E.2.1}
\]
The inverse image of the marked oval graph is the corresponding
finite cyclic graph covering, including flags and the selected
sheet.  Elements of $H$ give deck transformations of this
cover because they commute with $s$ and normalize $N$.

The graph construction is functorial: both cusp flags map by
the local section maps, interval classes by composition of
sections, and all incidence maps commute.  A peripheral
extension isomorphism therefore gives compatible isomorphisms
of these graph coverings and of their deck actions.  The base
graph isomorphism has degree $\epsilon=\pm1$.  Lifting one
circuit at each finite level gives
\[
 b(a)\equiv(a')^\epsilon\pmod{H'\cap b(N)}.
\]
The sign is the same at every level because the graph coverings
lie over that same base graph map.  The subgroups $H\cap N$
are cofinal in the full profinite topology of $H$, so
$b(a)=(a')^\epsilon$ after the conjugations identifying the
chosen sections.

For a dividing surface orient all ovals as the boundary of
one half.  The integral relative homology sequence gives
\[
 \ker\left(\widehat{\mathbb Z}^{,r}
       \longrightarrow H_1(S_g,\widehat{\mathbb Z})\right)
       =\widehat{\mathbb Z}(1,\ldots,1).                     \tag{E.2.2}
\]
There are two complementary halves; their boundary vectors
are the displayed diagonal vector and its negative, which
proves that these are all relations.  An isomorphism sends
each oval generator to a unit multiple of the target oval
generator.  Preservation of (E.2.2) forces all these unit
multipliers to agree.  On every marked oval the multiplier is
the already determined sign $\epsilon$, so the marked signs
agree.  Nondividing quotient surfaces allow the independent
boundary reversals.  For unordered markings the nonreal
peripheral classes occur in two-element conjugation orbits;
these record the number of interior marks of the quotient.
The quotient surface with these interior marks and with the
recorded boundary cyclic incidence determines the topological
marked type.  Doubling its homeomorphisms gives the desired
classification upstairs.  This proves both marked injections.
\end{proof}

\statement{E.3}{Proved}{The exact nonabelian-cohomology fibre of LM6 C.15}
Let $1\to K\to P^{\pm}\xrightarrow{q}Q^{\pm}\to1$ be a
continuous exact sequence of profinite groups with compatible
orientation characters, with $K\leq P$ their positive kernel.
Fix an orientation-reversing involution $\tau\in P^{\pm}$ and
put $\alpha=\operatorname{Ad}(\tau)$ and $\bar\tau=q(\tau)$.
For $k\in K$ satisfying $k\alpha(k)=1$, identify cocycles by
$k\sim b k\alpha(b)^{-1}$, $b\in K$.  Then
\begin{align}
 q(g)\star[k]&=[gk\alpha(g)^{-1}],\notag\\
 \{\text{$P$-classes over }[\bar\tau]\}
      &\cong C_Q(\bar\tau)\backslash H^1(C_2,K),             \tag{E.3.1}\\
 1\longrightarrow K^{k\tau}\longrightarrow C_P(k\tau)
      &\xrightarrow{q}
       \operatorname{Stab}_{C_Q(\bar\tau)}([k])
       \longrightarrow1 .                                  \tag{E.3.2}
\end{align}
Here $g\in P$ is any lift of the indicated quotient-centralizer
element.  The action need not fix the class of the trivial cocycle.

\begin{proof}
The fibre over the element $\bar\tau$ consists precisely of
$k\tau$ with $k\alpha(k)=1$.  Conjugation by $b\in K$ gives
the indicated equivalence.  If $q(g)$ centralizes $\bar\tau$,
then $gk\alpha(g)^{-1}\in K$, and
\[
             g(k\tau)g^{-1}=gk\alpha(g)^{-1}\tau.
\]
This proves the cocycle condition and the action law.  Replacing
$g$ by $bg$ changes the cocycle by the stated equivalence, so
the action is independent of lifts.  An involution whose image
is conjugate to $\bar\tau$ can first be conjugated by a lift
in $P$ to lie over that element.  Two such representatives
are $P$-conjugate exactly when their $K$-classes are in the
same quotient-centralizer orbit.  This proves (E.3.1).

The kernel of the centralizer map is $K^{k\tau}$.  Its image
fixes $[k]$.  If a quotient-centralizer element fixes $[k]$,
choose its lift $g$.  There is $b\in K$ with
$g(k\tau)g^{-1}=b(k\tau)b^{-1}$, and $b^{-1}g$ is a
centralizing lift.  This proves (E.3.2), including surjectivity
onto the stated stabilizer.

For the last qualification take
$P^{\pm}=D_8=\langle r,t\mid r^4=t^2=1,trt=r^{-1}\rangle$,
$P=\langle r\rangle$, $K=\langle r^2\rangle$, and $\tau=t$.
The two classes in $H^1(C_2,K)$ are $[1],[r^2]$.  Although
all of $Q=P/K$ centralizes $\bar\tau$, its element $q(r)$
interchanges these classes because $r\alpha(r)^{-1}=r^2$.
Thus the orbit set has one element and $C_P(t)=K$ maps
trivially to $Q$.  The attached finite script checks these
group operations and orbit sizes.
\end{proof}

\statement{E.4}{Proved}{The hyperelliptic comparison of LM6 C.16 uses only the element theorem}
Fix a hyperelliptic involution $\iota\in\Gamma_3$, and put
$H^{\pm}=C_{\Gamma_3^{\pm}}(\iota)$ and $H=H^{\pm}\cap\Gamma_3$.
The natural map $\widehat{H^{\pm}}\to P_3^{\pm}$ is injective,
with image $\overline{H^{\pm}}$.  Each orientation-reversing
involution in this image is $\overline H$-conjugate to a
discrete one.  For such a discrete involution $\tau$,
\[
 C_{P_3}(\tau)\cap\overline H
       =\overline{C_H(\tau)}.                               \tag{E.4.1}
\]
Each of the six types in (E.1.1) has a representative in $H^{\pm}$.
No centralizer equality for a nonsolvable finite subgroup is
used in these assertions.

\begin{proof}
Theorem~7.1 of \cite{CEBoggiHyperelliptic2024} states the
congruence subgroup theorem with the topology induced from
the action on the completed surface group.  Thus the
composite $\widehat H\to P_3\to Q_3$ is injective.  The first
map is injective, and adjoining the orientation quotient of
order two proves injectivity for $H^{\pm}$ as well.

We check the group to which the torsion theorem is applied.
The branched double cover gives the Birman--Hilden exact
sequence
\[
 1\longrightarrow\langle\iota\rangle\longrightarrow H
       \longrightarrow\Gamma_{0,[8]}\longrightarrow1.
                                                               \tag{E.4.2}
\]
The kernel is the deck group.  All sphere mapping classes
lift because the double-cover character sends each of the
eight meridians to the nonidentity element of $C_2$.
This is the exact sequence (3) in the same primary source.
Intersect the inverse image of the pure sphere group with
the level-three homology kernel.  The resulting finite-index
subgroup misses $\iota$, which acts as $-I$ on
$H_1(S_3,\mathbb F_3)$, and embeds as a finite-index subgroup
of the pure sphere mapping class group.

Here is a direct goodness verification.  After fixing three
points at $0,1,\infty$, the forgetful configuration fibration
has finitely generated free fundamental group as kernel and
a section, obtained by placing the new real point at
$2+\sum_j|z_j|$.  A finite-index subgroup of a finitely
generated free kernel contains a characteristic finite-index
subgroup.  The action on its finite quotient has finite image;
the resulting finite semidirect-product quotient of the whole
group detects that kernel quotient.  Thus completion of the
split extension is exact.  Induction and Hochschild--Serre
comparison, starting from free groups, prove goodness of
the pure sphere group.  They also prove finite cohomology
with finite coefficients and finite virtual cohomological
dimension.  Shapiro's lemma passes these properties to
finite-index subgroups, and Hochschild--Serre passes them to
finite extensions.  They therefore hold for $H^{\pm}$.
Filtered unions again give the comparison for arbitrary
discrete continuous $\mathbb F_2$-modules.

Theorem~A of \cite{CEBoggiZalesskii2025}, at $p=2$, now
gives torsion conjugacy and centralizer equality inside
$\widehat{H^{\pm}}$ for the one element $\tau$.  A negative
conjugator is made positive by multiplying by the source
involution.  Intersecting its centralizer equality with the
open positive subgroup gives (E.4.1).

For existence of all types use the smooth projective models
$y^2=f(x)$ with
\[
\begin{array}{c|c}
 (r,a)&f(x)\\\hline
 (0,1)&-(x^8+1)\\
 (1,1)&(x^2-1)(x^6+1)\\
 (2,1)&(x^2-1)(x^2-4)(x^4+1)\\
 (3,1)&(x^2-1)(x^2-4)(x^2-9)(x^2+1)\\
 (2,0)&x^8+1\\
 (4,0)&\displaystyle\prod_{j=1}^4(x^2-j^2).
\end{array}                                                   \tag{E.4.3}
\]
All polynomials are square-free of degree eight.  If there
are $2k>0$ real roots and positive leading coefficient, the
$k$ positive intervals on $\mathbb P^1(\mathbb R)$ lift to
$k$ ovals.  For $k<4$ the inverse image of each open
half-plane is connected because it contains branch points;
a negative real interval joins these halves through nonreal
points of the curve.  The complement of the ovals is
therefore connected.  The four-oval case is dividing by the
quotient Euler characteristic in E.1.  The negative polynomial
has no real points, including at infinity.  For $x^8+1$, the
two signs give two real circles, as the infinity coordinate
$y/x^4$ preserves each sign.  All points above the real
projective line are real, and the two connected half-plane
preimages are the two complementary halves.  Conjugation
commutes with the hyperelliptic deck involution in all six
models.  Every genus-three hyperelliptic involution is
conjugate to $\iota$, by the topological double-cover
construction with eight branch points, which proves the
claim for the specified $H^{\pm}$.
\end{proof}

\statement{E.5}{Proved}{The five oval centralizers of LM6 C.17}
Let $\tau\in\Gamma_3^{\pm}$ be a real involution with $r>0$,
and put $C=C_{\Gamma_3}(\tau)$.  The restriction of $q$ is an
isomorphism
\[
 \overline C^{\,P_3}\xrightarrow{\sim}
 C_{Q_3}(\bar\tau)=\overline C^{\,Q_3},\qquad
 C_{P_3}(\tau)=K_c^\tau\rtimes\overline C^{\,P_3}.           \tag{E.5.1}
\]
For these five types the quotient-centralizer action of E.3
fixes the trivial cohomology class.  Formula (E.5.1) does not
compute $K_c^\tau$.

\begin{proof}
Let $A$ be the oval multicurve.  Its curves are essential and
pairwise nonisotopic.  A contractible oval would bound a disc
which is exchanged with its other side by the involution,
forcing genus zero.  Parallel ovals would similarly force
the complementary component on either side to be an annulus,
forcing genus one.  Neither can occur here.

First take $b\in C_{Q_3}(\bar\tau)$.  It permutes the section
classes of the real surface extension and their centralizers.
To recover actual curve generators, choose a real marked point
on an oval and use the pointed comparison
\[
 \check\Gamma_{3,1}^{\pm}\cong
 \operatorname{Aut}(\Pi)\mathbin{\times}_{\operatorname{Out}(\Pi)}
                                      Q_3^{\pm}.             \tag{E.5.2}
\]
Corollaries~4.5 and 4.7 of \cite{CEBoggiCongruence2020}, for
the class of all finite groups, give the faithful pointed
action and its exact Birman sequence.  Its kernel is
$\Pi\cong\operatorname{Inn}(\Pi)$.  Lifting an outer element
and then adjusting its inner automorphism proves surjectivity
onto the displayed pullback, so these results give exactly
(E.5.2).  Adjoining the actual pointed reflection gives the
same identity for the orientation extension.

Lift $b$ in (E.5.2).  Its conjugate of the pointed reflection
is a section belonging to one of the actual target ovals,
by E.1.  Adjusting by an element of $\Pi$ makes these actual
pointed reflections agree.  The faithfulness in (E.5.2) makes
this an equality in the punctured procongruence group.
Applying E.2 with the one real marked point shows that $b$
takes the original oval generator to the target generator
or its inverse.  Repeating for every oval proves
$b\in(Q_3)_A$, the stabilizer of this actual unoriented
multicurve in the procongruence curve complex.

Put $L=\operatorname{Stab}_{\Gamma_3}(A)$, and let $L_0$ fix
each curve and its orientation.  Cutting and replacing the
boundary curves by labelled punctures gives the exact sequence
\[
 1\longrightarrow T\cong\mathbb Z^r\longrightarrow L_0
       \longrightarrow\prod_j\Gamma(S_j)^{\mathrm{pure}}
       \longrightarrow1,                                    \tag{E.5.3}
\]
where $T$ consists of the individual oval Dehn twists.  The
surface types are
\[
\begin{array}{c|c}
 (r,a)&\coprod_j S_j\\\hline
 (1,1)&S_{2,2}\\
 (2,1)&S_{1,4}\\
 (3,1)&S_{0,6}\\
 (2,0)&S_{1,2}\amalg S_{1,2}\\
 (4,0)&S_{0,4}\amalg S_{0,4}.
\end{array}                                                   \tag{E.5.4}
\]
Theorem~4.5 of \cite{CEBoggiProcongruence2014} identifies the
procongruence stabilizer with the closure of $L$ and gives
the exact completed version of (E.5.3), whose kernel is the
injected full $\widehat{\mathbb Z}^{,r}$.  Each component
has genus at most two.  Theorem~1.4(ii) of
\cite{CEBoggiCongruence2020}, with all finite groups as
formation, identifies its procongruence completion with the
full completion; the genus-zero comparison also follows
from Theorem~7.1 of \cite{CEBoggiHyperelliptic2024}.
Completion of (E.5.3) is right exact.  Its left group maps
injectively to the twist kernel just specified, so it is
also left exact.  The comparison of these two exact
sequences is an isomorphism at both ends and therefore in
the middle.  Accounting for the finite permutation and
orientation quotient gives
\[
                     \widehat L\cong(Q_3)_A.                 \tag{E.5.5}
\]
It also proves that $\widehat L\to P_3$ is injective, since
its composite with $q$ is injective.

We verify the torsion theorem hypotheses for this stabilizer,
not for the ambient genus-three group.  Pure sphere mapping
class groups are good by the proof in E.4; the genus-one
one-puncture group is virtually free; the genus-two closed
group is hyperelliptic, so the same finite-index sphere
argument applies.  Adding punctures preserves goodness by
the exact Birman sequence, goodness of surface groups, and
the Hochschild--Serre comparison.  The same proof gives
finite virtual cohomological type.  Hence all the component
groups in (E.5.3) are good, as are their product and the
twist group.  The exact completion already proved for
(E.5.3) gives goodness of $L_0$.  Finite extensions give
goodness and finite virtual two-cohomological type for
$\langle L,\tau\rangle$.

Theorem~A of \cite{CEBoggiZalesskii2025} now identifies the
centralizer of the one element $\tau$ in $\widehat L$ with
the closure of its discrete centralizer.  A discrete
centralizer preserves $A$, so its centralizer in $L$ is
exactly $C$.  Together with containment in $(Q_3)_A$ and
(E.5.5), this proves the quotient centralizer equality.
The copy of $\overline C$ in $P_3$ maps isomorphically onto
it.  The kernel of the full centralizer map is $K_c^\tau$;
the specified copy is a complement, proving (E.5.1).
Thus only the prime-two element case of the cited theorem
has been used.
\end{proof}

\statement{E.6}{Proved}{The free real type has the exact comparison of LM6 C.18}
For the fixed-point-free genus-three involution $\tau$, put
$N_4=S_3/\tau$, $E=\widehat{\pi_1(N_4)}$, and
$\Pi=\widehat{\pi_1(S_3)}$.  Let $\operatorname{Aut}_{Q}(E;\Pi)$
be the continuous automorphisms preserving $\Pi$ whose
restrictions have outer class in $Q_3^{\pm}$.  Then
\[
 \mathcal A:=\operatorname{Aut}_{Q}(E;\Pi)/\operatorname{Inn}(E)
   \cong C_{Q_3^{\pm}}(\bar\tau)/\langle\bar\tau\rangle
   \cong C_{Q_3}(\bar\tau).                                  \tag{E.6.1}
\]
Under this homeomorphism the closure of
$\operatorname{Mod}(N_4)$ equals the closure of
$C_{\Gamma_3}(\tau)$ in $Q_3$.  Equality of this closure
with $\mathcal A$ remains a density question.

\begin{proof}
The orientation-cover inclusion has full induced profinite
topology, as it has index two, so
$1\to\Pi\to E\to C_2\to1$ is exact.  Choose
$e\in E\setminus\Pi$, put $a=\operatorname{Ad}(e)|_\Pi$
and $z=e^2$.  Then $a^2=\operatorname{Inn}(z)$.  If
$b\in\operatorname{Aut}(\Pi)$ has outer class commuting
with $\bar\tau$, there is a unique $d\in\Pi$ with
$bab^{-1}=\operatorname{Inn}(d)a$, since $Z(\Pi)=1$.
Squaring gives $b(z)=d\,a(d)z$.  These two identities
show that $b$ extends to $E$ by $e\mapsto de$; they
preserve conjugation on $\Pi$ and the square of $e$.
The map is continuous and bijective on the two cosets
$\Pi$ and $\Pi e$.  It is unique by the same centre-free
argument.  Thus restriction gives the inverse image of
the outer centralizer.  Quotienting by inner automorphisms
of $\Pi$ removes $\operatorname{Inn}(\Pi)$, and inner
conjugation by $e$ removes the remaining class $\bar\tau$.
Finally
$C_{Q_3^{\pm}}(\bar\tau)=C_{Q_3}(\bar\tau)\times
\langle\bar\tau\rangle$, by the orientation character.
These constructions give (E.6.1); all continuous bijections
are homeomorphisms by compactness.

Every homeomorphism of $N_4$ has exactly one
orientation-preserving lift to $S_3$.  The lift is compatible
with composition and isotopy and identifies
$\operatorname{Mod}(N_4)$ with $C_{\Gamma_3}(\tau)$.
The precise mapping-class comparison, including its kernel
before choosing the positive lift, is recorded on page 1030
of \cite{CEStukow2016}.  Its completed action is exactly
the restriction in (E.6.1).  Taking closures proves the
last assertion without asserting the density itself.
\end{proof}

\statement{E.7}{Proved}{The obstruction to a hyperelliptic reduction in LM6 C.19}
For the positive lift
$j:\operatorname{Mod}(N_4)\cong C_{\Gamma_3}(\tau)$ and
every finite-index subgroup $H\leq\operatorname{Mod}(N_4)$,
the centralizer $C_{\Gamma_3}(j(H))$ has no nonidentity
finite-order element.  In particular $j(H)$ is never
contained in a hyperelliptic centralizer, and the standard
hyperelliptic subgroup of $\operatorname{Mod}(N_4)$ has
infinite index.

\begin{proof}
Glue two pairs of pants $P,Q$ along boundary circles
$a,b,c$, with orientation voltages $0,0,1$ on the three
edges of the resulting theta graph.  The orientation
double is connected, since the $a,c$ cycle has nonzero
voltage.  The base is a closed nonorientable surface of
Euler characteristic $-2$, hence $N_4$, and its double
is $S_3$.  The lifted pants graph has edges
\[
\begin{array}{c|c}
 a_+,b_+&P_+,Q_+\\
 a_-,b_-&P_-,Q_-\\
 c_+&P_+,Q_-\\
 c_-&P_-,Q_+.
\end{array}                                                   \tag{E.7.1}
\]
The six curves form a pants decomposition.  They are
essential and pairwise nonisotopic: a disc or annulus
bounded by them would be a union of complementary pairs
of pants, contradicting its Euler characteristic.

Choose orientations of the three base annuli and denote
their twists by $t_a,t_b,t_c$.  Their positive lifts are
\[
 j(t_a)=T_{a_+}T_{a_-}^{-1},\quad
 j(t_b)=T_{b_+}T_{b_-}^{-1},\quad
 j(t_c)=T_{c_+}T_{c_-}^{-1}.                                  \tag{E.7.2}
\]
The opposite signs arise because the deck involution
exchanges the two annuli and reverses the surface
orientation.  A finite-index $H$ contains positive powers
of all three twists.  If a positive finite-order mapping
class $f$ centralizes their lifts, uniqueness of signed
multitwist supports and exponents forces it to fix each
of the six individual curves.  The two exponents $m$ and
$-m$ cannot be interchanged by a positive mapping class.
This is precisely Lemma~3.17 of \cite{CEFarbMargalit2011}.

Each vertex in (E.7.1) has a distinct set of incident
labelled edges, so $f$ also fixes every complementary
pair of pants.  Its restrictions, with boundaries treated
as punctures, lie in the trivial pure three-punctured-sphere
mapping class groups.  By the cutting homomorphism,
Proposition~3.20 of the same source, $f$ lies in the
multitwist kernel $\mathbb Z^6$.  This group is torsion-free,
so $f=1$.  Proposition~2.3 there supplies the asserted
triviality for the pure three-punctured sphere.  The
standard nonorientable hyperelliptic subgroup lifts into
a hyperelliptic centralizer by the explicit maps on page
1030 of \cite{CEStukow2016}; it therefore has infinite
index.  The graph and the signed exponent vectors are
independently checked by \texttt{code/check\_real\_models.py}.
\end{proof}

\statement{E.8}{Proved}{Finite quotients retain the entire fixed group and nonabelian cohomology}
Let a finite group $C_2=\langle\tau\rangle$ act continuously
on a profinite group $K$, and let $N$ range over its open
normal $\tau$-invariant subgroups.  There are natural
homeomorphisms
\[
 K^\tau\cong\varprojlim_N(K/N)^\tau,
 \qquad
 H^1(C_2,K)\cong\varprojlim_N H^1(C_2,K/N).                  \tag{E.8.1}
\]
The cohomology set is given the quotient topology of its
compact cocycle space.  Neither the maps to individual
finite fixed groups nor the finite cohomology transition
maps are asserted to be surjective.

\begin{proof}
The fixed-group identity follows by applying the defining
equation $\tau(k)=k$ in every finite quotient.  For
cohomology, a compatible family of finite cohomology
classes defines, for each $N$, a nonempty closed subset
of $K$ consisting of all lifts of its cocycle class in
$K/N$.  Any finitely many such conditions are satisfied
by a lift from one common refinement, so compactness
gives $k\in K$ satisfying them all.  Its cocycle equation
holds modulo every $N$ and hence in $K$ itself.  This
proves surjectivity.  If two actual cocycles have the same
finite class at every $N$, the sets of conjugators
$b\in K$ satisfying $k'=bk\tau(b)^{-1}\pmod N$ are
nonempty closed sets with the same finite-intersection
property.  Their intersection gives an actual
conjugator, proving injectivity.  The orbit relation
is closed, as it is the image of a compact set, so
the cohomology space is compact Hausdorff.  The continuous
bijection to the inverse limit is therefore a
homeomorphism.
\end{proof}

\statement{E.9}{Proved}{The fixed congruence kernel is the fibre of oval stabilizers}
For one of the five oval types in E.5, let $A$ be its oval
multicurve and $L=\operatorname{Stab}_{\Gamma_3}(A)$.
Use E.5 to regard $\widehat L$ as a closed subgroup of
$P_3$ mapping isomorphically onto $(Q_3)_A$.  Then the
fibre over the identity coset of the continuous map
\[
 P_3/\widehat L\longrightarrow Q_3/(Q_3)_A
\]
is naturally the profinite set $K_c$.  Its $\tau$-fixed
subset is naturally $K_c^\tau$.

\begin{proof}
The inverse image of $(Q_3)_A$ is
$K_c\widehat L$, and $K_c\cap\widehat L=1$.  Every coset
in this fibre therefore has a unique representative
$k\in K_c$.  This is a homeomorphism by compactness.
The actual involution $\tau$ normalizes $L$ and its
completion.  Its action on the fibre takes $k\widehat L$
to $\tau(k)\widehat L$.  This equals $k\widehat L$
exactly when $k^{-1}\tau(k)\in K_c\cap\widehat L=1$.
This proves the fixed-subset assertion.  In particular,
a theorem saying that there is just one fixed lift of
this geometric multicurve would prove $K_c^\tau=1$;
preservation of the procongruence multicurve alone does
not say that.
\end{proof}

\statement{E.10}{Open}{The six congruence-kernel computations in prompt E.1}
For each type (E.1.1), determine
$K_c^\tau$ and the nonabelian pointed set
$H^1(C_2,K_c)$ for $\operatorname{Ad}(\tau)$.
E.3 gives the exact class and centralizer formula after
these groups and sets are known; E.5 supplies its
quotient-centralizer term for the five oval types.
For $(0,1)$ the density question in E.6 remains an
additional unresolved term.  The current state is
\[
\begin{array}{c|c|c|c}
 (r,a)&C_{Q_3}(\bar\tau)&K_c^\tau&H^1(C_2,K_c)\\\hline
 (0,1)&\mathcal A\text{ of E.6}&\text{uncomputed}&\text{uncomputed}\\
 (1,1)&\overline C^{\,Q_3}&\text{uncomputed}&\text{uncomputed}\\
 (2,1)&\overline C^{\,Q_3}&\text{uncomputed}&\text{uncomputed}\\
 (3,1)&\overline C^{\,Q_3}&\text{uncomputed}&\text{uncomputed}\\
 (2,0)&\overline C^{\,Q_3}&\text{uncomputed}&\text{uncomputed}\\
 (4,0)&\overline C^{\,Q_3}&\text{uncomputed}&\text{uncomputed}.
\end{array}                                                   \tag{E.10.1}
\]
The existence of six geometric classes is proved, and
their distinctness persists in both completions; a
classification of all profinite classes does not follow.
E.8 gives exact inverse-limit computations but no
presentation of the required finite quotients of $K_c$.
E.9 identifies a separate geometric fixed-fibre problem.
The hyperelliptic route settles the intersection in
E.4; E.7 rules out extending that route to the free
type by passing to a finite-index nonorientable
mapping class subgroup.  No vanishing of the genus-three
congruence kernel is assumed.  The contemporary source
\cite[Problem~19 and the following paragraph]{CEAndrewHenselHughesWade2026}
explicitly records that the mapping class congruence
subgroup problem remains open in genus at least three.


% E.30 onward reserved for the five-point census.
\subsection{Five marked points and cote 148}
\label{sec:E2}

We distinguish three classification problems.  The first concerns
a complex sphere with a five-point set $K$.  The second equips
that configuration with an antiholomorphic involution.  The third
concerns finite subgroups of the extended mapping class group of
the five-punctured sphere.  These are different problems.
An antiholomorphic map of finite order can have order greater than two.
We write $D_n$ for the
dihedral group of order $2n$ and $\mu_n$ for the complex $n$th roots
of unity.  The manuscript passages here are cote 148, archival
pp.~30--35, scans 31--36; its separate classification occupies
handwritten pp.~1--3, scans 34--36.  We record below which conclusions
come from independent mathematics and which diagram entries require
interpretation of the manuscript's additional decorations.  We inspected
the original Montpellier PDF directly \cite{DE_Grothendieck148}.  Its printed page numbers are
one less than the one-based PDF scan numbers; the supplied first-pass
transcription shifts the archival page numbers down by one further unit.

\statement{E.30}{Proved}{The holomorphic classification of five-point configurations}
The possible full holomorphic automorphism groups are
\[
 1,\quad C_2,\quad C_4,\quad D_3,\quad D_5.
\]
If the maximum order $\nu$ of a nonidentity holomorphic automorphism
is $5$, $4$, or $3$, there is exactly one isomorphism class, respectively
\[
 \mu_5,\qquad\{\infty\}\cup\mu_4,
 \qquad\{0,\infty\}\cup\mu_3.
\]
Their full groups are $D_5,C_4,D_3$.  The labeled structures are,
respectively, an unoriented pentagon, a distinguished point with an
oriented square on the other four points, and a three-set together
with its two-element orientation torsor.  These assertions verify
the holomorphic classification on cote 148, scan 34, handwritten p.~1.
\begin{proof}
The holomorphic group acts faithfully on $K$, since a M\"obius
transformation fixing three points is the identity.  It is finite.
A nonidentity finite-order M\"obius map is conjugate to
$z\mapsto\zeta_\nu z$.  Its only fixed points are $0,\infty$, and
all other orbits have length $\nu$.  Hence
$5=f+m\nu$, with $0\leq f\leq2$, $m\geq1$, giving
$\nu\in\{2,3,4,5\}$ and the three displayed normal forms for
$\nu\geq3$.

For $K=\mu_5$, every automorphism preserves its unique generalized
circle and hence its cyclic order up to reversal.  It acts on the
five points through $D_5$.  The maps $z\mapsto\zeta_5z$ and
$z\mapsto1/z$ realize all of $D_5$.
For $K=\{\infty\}\cup\mu_4$, the unit circle is the unique
generalized circle containing four marked points: any other
four-set contains $\infty$, and its other three points are not
collinear.  Thus $\infty$ is invariant.  An automorphism is affine;
preservation of the centroid of $\mu_4$ forces its translation
part to vanish, leaving exactly multiplication by $\mu_4$.
For $K=\{0,\infty\}\cup\mu_3$, deleting either pole gives an
equianharmonic four-set, with $j=0$, whereas deleting a root gives
$j=2048/3$.  This follows by substitution in
\[
 j(\lambda)=256\frac{(1-\lambda+\lambda^2)^3}
 {\lambda^2(1-\lambda)^2}.
\]
Therefore the pole pair is intrinsic.  An automorphism has form
$az$ or $a/z$, and preservation of $\mu_3$ forces $a\in\mu_3$.
The resulting group is $D_3$.

If every nonidentity element has order two, the group is elementary
abelian.  After one involution is written as $z\mapsto-z$, every
commuting involution is either that map or $z\mapsto a/z$.
There are at most four elements.  A Klein four action has no common
fixed point and all its point orbits have even cardinality, so cannot
preserve a five-set.  Thus the remaining nontrivial full group is
$C_2$.  The stated labeled structures have stabilizers respectively
$D_5,C_4,S_3$ in $S_5$, with the same actions as the normal forms;
this also proves the categorical identifications used in scan 34.
\end{proof}

\statement{E.31}{Proved}{The one-parameter involution locus and the square and tetrahedron cases}
Every five-point configuration with a nontrivial holomorphic
involution can be written
\[
 K_t=\{\infty,1,-1,t,-t\},\qquad t\ne0,\ \ t^2\ne1.
\]
For a configuration with the involution specified, its isomorphism
class is determined by
\[
 \lambda=t^2,\qquad\beta=\lambda+\lambda^{-1}\in\mathbb C\setminus\{2\}.
\]
The full group is $C_2$ except at
\[
 \begin{array}{c|ccc}
 \beta&-2&-10/3&18\\\hline
 \operatorname{Aut}^+(K_t)&C_4&D_3&D_5.
 \end{array}
\]
A harmonic four-set together with a fixed point of one of its
fixed-point-free involutions gives $\beta=-2$ or $34$.
The second case has full group $C_2$; it is neither the pentagonal
nor the square-pyramidal case.  An equianharmonic four-set with
such a fixed point gives $\beta=-14$ and full group $C_2$.
These formulas complete the calculation begun and then crossed out
on cote 148, scan 35, handwritten p.~2, and answer its question
about adding an edge midpoint to a spherical square.
\begin{proof}
An involution preserving five marks has exactly one marked fixed
point.  Send it to $\infty$ and the other fixed point to $0$,
write the involution as $z\mapsto-z$, and scale one two-orbit to
$\{1,-1\}$.  Changing the chosen two-orbit replaces $t$ by
$\pm t^{-1}$; changing its ordering replaces $t$ by $-t$.
Conversely any isomorphism preserving the specified involution
must preserve its marked and unmarked fixed points individually,
so is a scaling.  These are exactly the identifications.  Hence
$\beta$ is complete and $\beta=2$ is exactly the forbidden collision.
The normal forms of E.30 give $t=i$, $i/\sqrt3$, and $\sqrt5-2$,
respectively, yielding the three exceptional values.  E.30 excludes
any other enhancement.  All involutions in each of $C_4,D_3,D_5$
are conjugate within that group, so forgetting the chosen involution
does not introduce another value of $\beta$.

For the four-set $\{1,-1,t,-t\}$ a cross-ratio is
$((1-t)/(1+t))^2$, and direct substitution gives
\[
 j=16\frac{(\lambda^2+14\lambda+1)^3}
 {\lambda(\lambda-1)^4}.
\]
The equianharmonic condition is $\lambda^2+14\lambda+1=0$,
namely $\beta=-14$.  The harmonic condition factors exactly as
\begin{align*}
 &16(\lambda^2+14\lambda+1)^3
 -1728\lambda(\lambda-1)^4\\
 &\hspace{15mm}=16(\lambda+1)^2(\lambda^2-34\lambda+1)^2,
\end{align*}
giving $\beta=-2$ or $34$.  In the square model $\mu_4$, the
involution $-z$ gives the pole and $\beta=-2$; an involution such
as $i/z$ gives a point halfway between consecutive vertices and
$\beta=34$.  For the tetrahedral four-set the three involutions
have six fixed points, the vertices of the associated octahedron;
adding any one gives $\beta=-14$.
The exact symbolic program \texttt{check\_special\_configurations.py}
enumerates every M\"obius candidate from its images on three marks
and verifies all six listed normal forms, including the generic
control $t=2$, with no numerical tolerance.
\end{proof}

\statement{E.32}{Refuted}{Every finite antiholomorphic automorphism preserving five marks has order two}
The inference in cote 148, scan 35, handwritten p.~2, that an
even-order antiholomorphic automorphism preserving five points
must fix a marked point and therefore have order two is false.
The exact possible orders of antiholomorphic automorphisms of
five-point configurations are $2,6,10$.
\begin{proof}
Let $\omega=e^{2\pi i/3}$ and
$K=\{0,\infty,1,\omega,\omega^2\}$.  The map
\[
 \sigma(z)=\frac{\omega}{\overline z}
\]
interchanges $0$ and $\infty$ and cyclically permutes the three
roots.  Its marked permutation has cycle type $(2)(3)$;
$\sigma^2(z)=\omega^2z$, so its exact order is six.
It has no fixed point: a finite nonzero fixed point would satisfy
$|z|^2=\omega$, which is impossible, and the two poles are exchanged.
For $K=\mu_5$ the map
\[
 \tau(z)=\frac{\zeta_5}{\overline z}
\]
has order ten.  Its restriction to $K$ is a five-cycle, while
$\tau^5(z)=1/\overline z$ fixes all five marks.  Thus restriction
to the marks alone is not a faithful representation of the extended
automorphism group.  The pair consisting of that permutation and
the orientation sign is faithful: an element in its kernel is a
holomorphic map fixing five points.

For completeness, any antiholomorphic automorphism has even order,
and its square is in one of the positive groups in E.30.  Orders
four or eight would force a marked fixed point, since a permutation
of five letters whose order divides four and which has no fixed
point cannot occur.  A finite-order anti-M\"obius map with a fixed
point is an involution: send the point to $\infty$, write the map
as $a\overline z+b$, and observe that its square has positive real
linear coefficient $|a|^2$, so finite order forces the square to
be the identity.  The remaining possibilities are $2,6,10$, all
realized above or by conjugation.  The finite semilinear and
permutation computations are checked in \texttt{check\_census.py}.
\end{proof}

\statement{E.33}{Proved}{Real involutions, dimensions, and their exact relation to the symmetric square}
An antiholomorphic \emph{involution} preserving five marked points
has a fixed circle and fixes exactly $r=1,3,$ or $5$ of the marks.
Each of these three spaces of real pairs $(\mathbb P^1,K,\sigma)$
has real dimension two.  The examples in cote 148, scan 36,
handwritten p.~3, are valid: the pentagonal configuration has an
$r=5$ real structure, the triangular bipyramid has an $r=3$
structure, and the square pyramid has an $r=1$ structure.
These examples are not exclusive to the indicated type.
\begin{proof}
An involution of an odd set has a fixed point.  Hence $\sigma$ has
a fixed point on the sphere and is M\"obius-conjugate to complex
conjugation.  The marks off its fixed circle occur in conjugate
pairs, giving $r=1,3,5$.  For $r=5$, five distinct real points
modulo $\operatorname{PGL}_2(\mathbb R)$ have dimension $5-3=2$.
For $r=3$, choose and order the real marks as $0,1,\infty$; then
\[
 K=\{0,1,\infty,\lambda,\overline\lambda\},
 \qquad\operatorname{Im}\lambda>0.
\]
The remaining finite changes of normalization do not alter dimension.
For $r=1$, place the real mark at $\infty$; two distinct points
in the upper half-plane and their conjugates give four real
parameters, modulo the two-dimensional real affine group.
Again the finite exchanges do not alter dimension.

The complete real-structure counts on the three enhanced
configurations are
\[
\begin{array}{c|c|c}
 K&\operatorname{Aut}^{\pm}(K)&
 \text{fixed-mark counts of its negative involutions}\\\hline
 \mu_5&D_5\times C_2&5,1,1,1,1,1\\
 \{0,\infty\}\cup\mu_3&D_3\times C_2&3,3,3,3\\
 \{\infty\}\cup\mu_4&D_4&3,3,1,1
\end{array}
\]
For example $z\mapsto1/\overline z$ fixes $\mu_5$, conjugation
fixes the two poles and $1$ in the bipyramid, and
$z\mapsto i\overline z$ fixes only $\infty$ in the square pyramid.
The other entries follow by multiplying these maps by the positive
groups of E.30 and counting their fixed marks.

One must also qualify the real-locus shorthand on scan 36(a).
After three ordered real marks are normalized, the model
$\operatorname{Sym}^2_{\mathrm{dist}}(
\mathbb P^1\setminus\{0,1,\infty\})$ parametrizes an unordered
remaining pair.  Its real points include both two real points and
a nonreal conjugate pair.  Thus its entire real locus is not the
$r=5$ locus: the pair $\{i,-i\}$ is an explicit $r=3$ counterexample
to that unqualified identification.  The $r=5$ locus is its
split-real part.  Fully unordered five-point moduli additionally
identifies different choices of the normalized triple.
\end{proof}

\statement{E.34}{Proved}{The generic real involution locus}
For $K_t$ of E.31 with full positive group $C_2$, a real structure
exists exactly when $\beta$ is real.  The full extended group is
then $C_2\times C_2$, and the fixed-mark counts of its two negative
involutions are
\[
\begin{array}{c|c}
 \text{range of }\beta&\text{counts}\\\hline
 \beta>2&\{5,1\}\\
 \beta<-2&\{3,3\}\\
 -2<\beta<2&\{1,1\}.
\end{array}
\]
The enhanced values in E.31 are omitted from the indicated ranges.
\begin{proof}
The invariant of the conjugate configuration is $\overline\beta$,
and E.31 proves completeness of $\beta$.  Thus a negative
isomorphism implies $\beta=\overline\beta$.  Conversely if
$\beta>2$, take $t$ real and use $\overline z$ and
$-\overline z$.  If $\beta<-2$, take $t$ purely imaginary and
use the same two maps.  If $-2<\beta<2$, choose $|t|=1$ and use
$t\overline z$ and $-t\overline z$.  Their respective fixed-mark
counts are exactly the displayed values.  An anti-isomorphism
normalizes the unique positive involution; it fixes its unique
marked pole and therefore has the form $a\overline z$ in this
normalization.  It is an involution.  The negative coset has two
elements, proving both exhaustiveness and the group assertion.
\end{proof}

\statement{E.35}{Proved}{Reflection section classes for five punctures}
Let $U=\mathbb P^1\setminus K$ and let $\sigma$ be a real involution
fixing $r\in\{1,3,5\}$ marks.  Put $m=(5-r)/2$.  Its discrete
orbifold extension is
\[
 1\longrightarrow\pi_1(U)\longrightarrow
 E_\sigma\longrightarrow C_2\longrightarrow1,
 \qquad
 E_\sigma\cong\Bigl(\mathop{*}_{j=1}^{r}C_2\Bigr)*F_m.
\]
The sign map is nontrivial on each $C_2$ factor and trivial on
$F_m$.  It has exactly $r$ conjugacy classes of sections under
$\pi_1(U)$, one for each component of $U^\sigma$; every section
has trivial fixed subgroup in $\pi_1(U)$.
This verifies the five-point instances of LM4 B.4 and B.5.
\begin{proof}
The quotient of the sphere by the reflection is a disk.  The $r$
real punctures split its mirror boundary into $r$ open intervals;
the $m$ conjugate pairs give $m$ interior punctures.  Cut from each
interior puncture to the boundary, and use paths from one interior
basepoint to the $r$ mirror intervals.  Orbifold van Kampen gives
one reflection generator of order two per interval and one free
orientation-preserving generator per interior puncture, with no
other relation.  The index-two kernel is torsion-free, and its rank
is $(r-1)+2m=4$, as required for a five-punctured sphere.

Reduced-word normal form in a free product shows that every
finite-order element is conjugate into one finite factor.  Distinct
factors have nonconjugate involutions, and the centralizer of a
factor involution is that factor itself.  Each $E_\sigma$-conjugacy
class of such an involution is already one $\pi_1(U)$-conjugacy
class: for a sign-negative conjugator $g$ of $t_j$, write
$g=(gt_j)t_j$ with $gt_j\in\pi_1(U)$.  Therefore there are exactly
$r$ section classes, and intersecting its centralizer with the
kernel gives $1$.  The quotient construction identifies the
classes with the $r$ open fixed intervals, giving the claimed
geometric correspondence.
\end{proof}

\statement{E.36}{Proved}{The twenty finite-subgroup classes in type $(0,5)$}
The discrete extended mapping class group of a sphere with five
unordered punctures has exactly twenty conjugacy classes of finite
subgroups, including the trivial group.  Seven preserve orientation
and thirteen have nontrivial sign.  In the following table, the
last column lists the fixed-mark counts of all negative involutions,
with multiplicity; an empty entry means there are none.
\begin{center}\small
\begin{tabular}{c|r|c|l}
 $G$&$|G|$&$G^+$&\text{negative involutions}\\\hline
 $1$&1&$1$&\\
 $C_2$&2&$C_2$&\\
 $C_3$&3&$C_3$&\\
 $C_4$&4&$C_4$&\\
 $C_5$&5&$C_5$&\\
 $D_3$&6&$D_3$&\\
 $D_5$&10&$D_5$&\\\hline
 $C_2$&2&$1$&$5$\\
 $C_2$&2&$1$&$3$\\
 $C_2$&2&$1$&$1$\\
 $C_2\times C_2$&4&$C_2$&$5,1$\\
 $C_2\times C_2$&4&$C_2$&$3,3$\\
 $C_2\times C_2$&4&$C_2$&$1,1$\\
 $C_6$&6&$C_3$&$3$\\
 $D_3$&6&$C_3$&$3,3,3$\\
 $D_4$&8&$C_4$&$1,1,3,3$\\
 $C_{10}$&10&$C_5$&$5$\\
 $D_5$&10&$C_5$&$1,1,1,1,1$\\
 $D_3\times C_2$&12&$D_3$&$3,3,3,3$\\
 $D_5\times C_2$&20&$D_5$&$1,1,1,1,1,5$
\end{tabular}
\end{center}
Thus the type $(0,5)$ total in LM4 B.26 is correct, although it
must not be identified with the number of full automorphism groups
of individual complex configurations.
\begin{proof}
Kerckhoff's realization theorem allows orientation-reversing mapping
classes, and its extension to punctures is explicit in
\cite[Section IV, pp.~255--256; Theorem 5]{DE_Kerckhoff1983}.
Hence a finite subgroup is realized by holomorphic and antiholomorphic
automorphisms of a five-punctured sphere, and extends over its
punctures.  E.30 shows that its positive subgroup is one of
$1,C_2,C_3,C_4,C_5,D_3,D_5$.
This agrees with the positive conjugacy classification in
\cite[Theorems 3--4]{DE_Stukow2004}.  We give the remaining
classification directly.

If $G^+=1$, a negative group is generated by a reflection, with
$r=1,3,5$, giving three classes.  If $G^+=C_2$, its unique marked
fixed pole is preserved by all of $G$.  In coordinates in which
$G^+=\langle-z\rangle$, any negative normalizer preserving that
pole has form $a\overline z$ and has order two.  The full group
is a Klein four group with two perpendicular reflection circles.
The other four marks are either one free four-orbit, two two-orbits
on one mirror, or one two-orbit on each mirror.  These give exactly
$\{1,1\},\{5,1\},\{3,3\}$.

For $G^+=C_3$, the marks are the two poles and one three-orbit.
A negative normalizer either fixes the poles separately, giving
the dihedral group generated by $\zeta_3z$ and $\overline z$,
or exchanges them, giving the cyclic group generated by
$\zeta_3/\overline z$.  For $G^+=C_4$, the single marked pole
must be fixed, and the only extension is the dihedral group
$D_4$.  For $G^+=C_5$, the same normalizer calculation gives
$D_5$ or $C_{10}$, according as the two poles are fixed separately
or exchanged.  Finally the $D_3$ and $D_5$ positive actions have
the rigid configurations of E.30; their full extended groups are
$D_3\times C_2$ and $D_5\times C_2$.  Any negative extension
already has twice the positive order, so must be the full group.
This proves that the table is exhaustive.

These are conjugacy classes, not merely abstract groups.  For the
cyclic positive actions, mark orbits of equal kind move into each
other in the quotient sphere while the poles are retained.  For
the Klein four cases the quotient is a mirror lune: its marked
corner and the numbers of points on its two mirror edges and in
its interior determine the decoration up to homeomorphism.  The
three distributions just listed exhaust its possibilities.  The
other actions are rigid as described above.  A homeomorphism of
a decorated quotient lifts to a conjugacy of its sphere action.
Conversely the permutation action on the five punctures together
with the orientation sign is invariant under mapping-class
conjugacy.  The twenty models give distinct conjugacy classes of
subgroups of $S_5\times C_2$, as checked exactly in
\texttt{check\_census.py}; hence none of the twenty classes can
merge in the mapping class group.  The induced finite mapping-class
action is faithful, since a positive element in its kernel fixes
all five marks, and a negative element cannot be isotopic to the
identity.  This finishes the classification.
\end{proof}

\statement{E.37}{Proved}{The nine vertex families contain 542 labeled objects}
The vertex structures displayed in cote 148, scan 31, have the
following cardinalities and are each one transitive $S_5$-set:
\[
\begin{array}{c|r|r}
 \text{type}&\text{stabilizer order}&\text{cardinality}\\\hline
 I_{10}&10&12\\
 I_6&6&20\\
 I_4&4&30\\
 II_2&2&60\\
 II'_2&2&60\\
 II_1&1&120\\
 III_2&2&60\\
 III_1&1&120\\
 III'_2&2&60
\end{array}
\]
Their disjoint union has $542$ objects in nine $S_5$-orbits.
\begin{proof}
The first three counts can be read without a geometric drawing:
there are $(5-1)!/2=12$ unoriented pentagons; there are
$\binom53\,2=20$ choices of a triangle with its orientation
torsor identified with the remaining pair; and there are
$5\cdot(4-1)!=30$ distinguished-point oriented squares.
A pentagon with one distinguished vertex has $12\cdot5=60$
structures.  Equivalently choose a three-set $J$, a point $P\in J$,
and a bijection from the complementary two-set to $J\setminus P$:
$10\cdot3\cdot2=60$.  These are the parameters of $II_2,II'_2$
and $III_2$ in the diagram, with their respective geometric
realizations.  The fully pinned objects $II_1,III_1$ have trivial
stabilizer and therefore $5!=120$ labelings.  For $III'_2$, choose
$P$ in five ways and an oriented square with the extra codiagonal
choice on the remaining four points in twelve ways, giving $60$;
its stabilizer is the indicated double transposition.  Equivalently
all nine sets are the coset sets $S_5/H$ for the listed stabilizers.
Every set is transitive by relabeling, so orbit--stabilizer gives
the whole table.  Summing the three levels gives
$62+240+240=542$.  \texttt{check\_census.py} constructs and
checks these finite $S_5$-sets exactly.
\end{proof}

\statement{E.38}{Refuted}{The fourteen arrow families form fourteen $S_5$-orbits}
For the arrow families and local degrees displayed in cote 148,
scans 31 and 33, the sizes are
\[
 60,120;\quad60,120,60,60,120;\quad
 120,120,120,120,360,120,60.
\]
They sum to $1620$.  They cannot form fourteen $S_5$-orbits.
The size-$360$ family from $III_1$ to $II_1$ has three free
$S_5$-orbits, and the indicated remaining families give thirteen
orbits.  The resulting count is sixteen, with sizes
$5\cdot60+11\cdot120$.
\begin{proof}
The two $f$-families contribute $60+120=180$ arrows.  The five
middle families contribute $420$.  The seven lower families,
including the family with three arrows out of each of the
$120$ fully pinned $III_1$-objects, contribute $1020$.
This gives $180+420+1020=1620$.

An orbit of a group of order $120$ cannot have $360$ elements.
More precisely, stabilizers of arrows are contained in the
stabilizers of their sources.  Since $III_1$ is a free $S_5$-set,
every arrow with this source has trivial stabilizer.  Thus its
$360$ arrows split into exactly three orbits of size $120$.
For the remaining families, the indicated auxiliary datum is
transitive under the stabilizer of the source: it is a vertex of
a triangle, one of two equivalent choices, or a uniquely determined
choice, as appropriate.  Their arrow sets are correspondingly
$S_5/C_2$ or $S_5$, with the sizes listed above.  Therefore the
fourteen \emph{families} give $13+3=16$ orbits.

The computation file records each source, target, cardinality,
source degree and target degree, and constructs the corresponding
coset-orbit model.  This arithmetic does not establish relations
among the geometric arrows.  Direct inspection of
scan 33 confirms the total $1620$ and the $3\times120$ entry,
although overwritten intermediate products are inconsistent and hard
to read.  Crucially, the word \emph{orbites} following $14$ is
struck out in the original.  Thus the summary in R9, which presents
fourteen orbits as a clean assertion, is too strong.  Fourteen is
the correct number of displayed arrow families, and sixteen is the
correct number of $S_5$-orbits for that diagram.  We do not attribute
the uncancelled fourteen-orbit assertion to the manuscript.
\end{proof}

\statement{E.39}{Refuted}{The undecorated square pyramid has no real structure}
If the phrase on cote 148, scan 31, excluding $I_4$ from the real
structures is interpreted as a statement about the undecorated
five-point complex configuration, it is false.  For
$K=\{\infty,1,i,-1,-i\}$, complex conjugation preserves $K$,
and its full extended automorphism group is $D_4$ of order eight.
\begin{proof}
Conjugation fixes $\infty,1,-1$, exchanges $i,-i$, and normalizes
the positive generator $z\mapsto iz$ by inversion.  The group
it generates with that rotation is $D_4$.  Since E.30 gives
positive order four, there are no further negative elements.
Thus $G=\widetilde G=C_4$ is not correct for the undecorated
configuration when $\widetilde G$ means all holomorphic and
antiholomorphic automorphisms.  An additional oriented decoration
can exclude the reflection, in which case the statement concerns
a different object.  The scan-31 diagram uses oriented square
data, whereas scan 36 explicitly includes the square pyramid
among real configurations.  We retain this distinction rather
than identify two different categories.
\end{proof}


\statement{E.40}{Proved}{The stable-tree counts on scan 36}
For five labeled marks, the stable genus-zero dual trees with zero,
one, and two edges have respectively $1,10,15$ combinatorial types.
The corresponding decorated counts written at the bottom of cote 148,
scan 36, are $12$ in the pentagonal top type, $60$ and $120$ for the
two indicated one-edge decorations, and $60$ for the indicated
two-edge decoration.  The additional pyramid and bipyramid counts
are $30$ and $20$.
\begin{proof}
A one-edge stable tree splits the labels as $2+3$, giving
$\binom52=10$.  A two-edge tree has a single central mark and two
unordered pairs at its ends: $5\cdot3=15$.  There are no further
stable tree shapes for five marks.  The stabilizers of a fixed
one-edge and two-edge labeled partition have orders $12$ and $8$.
The decorated objects of E.37 have respective stabilizers $2,1$
above the one-edge type and $2$ above the indicated two-edge type.
Thus their fibers contain $12/2=6$, $12/1=12$, and $8/2=4$
objects; multiplication gives $60,120,60$.  The other three counts
are the first three entries of E.37.  This verifies the displayed
numerical statements without identifying the distinct decorated
families having the same cardinality.
\end{proof}

\statement{E.50}{Proved}{LM4 D.28 holds for the full pure three-cusp group}
Let \(p\) be a prime and \(\rho:C_p\to\mathcal P\) a
continuous homomorphism whose multiplier has
\(p\)-component one. Let \(\alpha\) be the unique normalized
actual representative of the image of a named generator,
and set
\[
 E=F\rtimes_\alpha C_p,\qquad
 S=\{1\}\times C_p,\qquad C=C_F(S).
\]
There is one \(F\)-conjugacy class of sections of
\(E\to C_p\). Each of the three peripheral normalizers
contains a section, and inclusion induces
\[
                     C^{(p)}\xrightarrow{\sim}F^{(p)}.
\]
Here \(C^{(p)}\) is the maximal pro-\(p\) quotient of
the actual full fixed subgroup. The statement is valid
for \(\mathcal P\), not merely for \(\widehat{GT}_0\).

\begin{proof}
E.51 proves the normalized homomorphic section without
using LM4 D.32. Hence \(\alpha^p=1\), and E.51's
Magnus argument makes \(\alpha\) identity on \(F^{(p)}\).
There is consequently a continuous map
\[
 E\longrightarrow F^{(p)}\times C_p,
 \qquad (g,t^j)\longmapsto(\bar g,t^j).
\]
Write \(c_1=x,c_2=y,c_3=(xy)^{-1}\) and
\(L_i=\overline{\langle c_i\rangle}\). By E.52,
\(N_F(L_i)=L_i\). Peripheral preservation thus gives
\(1\to L_i\to N_E(L_i)\to C_p\to1\).
Choose \(a_i\) over the named generator and write
\[
 a_i c_i a_i^{-1}=c_i^\lambda,
 \qquad a_i^p=c_i^{m_i},
 \qquad (\lambda-1)m_i=0.
\]
In the displayed pro-\(p\) quotient,
\(a_i\) has first coordinate \(\bar c_i^{v_i}\),
by the pro-\(p\) instance of E.52. Hence
\((m_i)_p=pv_i\). For
\(N_\lambda=1+\lambda+\cdots+\lambda^{p-1}\),
the equation \(N_\lambda u_i=-m_i\) is now solvable
in \(\widehat{\mathbb Z}\): at a prime \(q\) with
\(\lambda_q\ne1\) both the norm and \((m_i)_q\)
are zero; at every remaining prime the norm is \(p\),
and the only nonunit division was supplied above.
Thus \(a_ic_i^{u_i}\) has order \(p\) and defines
a peripheral section \(S_i\).

Let \(e_p\in\widehat{\mathbb Z}\) have coordinate
one at \(p\) and zero elsewhere. The group \(S_i\)
centralizes \(c_i^{e_p}\), whose mod-\(p\)
abelian image is the nonzero cusp vector.
Consequently \(H^1(C_F(S_i),\mathbb F_p)\) has
dimension at least one. For completeness, the
cohomological input used in the refereed proof is
Kochloukova--Zalesskii, Theorem~2 and Section~1.1,
equations (1)--(6), in the opened author manuscript.
It applies to full profinite groups with no
\(C_p^2\) and an open subgroup of finite
\(p\)-cohomological dimension.\cite{A_KZ2007}
Here every finite subgroup of \(E\) injects into
\(C_p\), and \(\operatorname{cd}_pF=1\).
For an order-\(p\) subgroup \(T\),
\(N_E(T)=C_F(T)\times T\). Restriction, K\"unneth,
and the Hochschild--Serre bound therefore give
\[
 \sum_{[T]}
 \bigl(1+\dim_{\mathbb F_p}H^1(C_F(T),\mathbb F_p)\bigr)
       =\dim_{\mathbb F_p}H^n(E,\mathbb F_p)\le3
       \qquad(n>1).
\]
The sum is finite: the restriction map is a dense
embedding from a finite-dimensional space into the
product, so it is surjective and that product is finite.

Two distinct cusp-section classes would contribute
at least four. They are therefore all conjugate.
Conjugating the \(x\)- and \(y\)-sections to one
representative places two elements with independent
mod-\(p\) abelian cusp vectors in its fixed subgroup.
Its contribution is at least three; the bound leaves
exactly this class and gives fixed-subgroup first
cohomology dimension two. Conjugacy under \(E\)
equals conjugacy under \(F\), since \(E=FT\).

The image of \(C\) in \(F^{(p)}\) is closed and
surjects onto its Frattini quotient, hence is all
of \(F^{(p)}\). Closed subgroups of the projective
group \(F\) are projective. The maximal pro-\(p\)
quotient of \(C\) is thus free pro-\(p\), with
rank two by the cohomology calculation. The induced
surjection between free pro-\(p\) groups of rank
two is an isomorphism by the Hopf property.
No \(S_3\)-symmetry equation or pentagon was used.
\end{proof}

The original LM4 D.28 register row restricted its
domain to \(\widehat{GT}_0\); E.50 records the domain
actually needed in LM4 D.32. Once R1 is established,
the displayed hypothesis forces \(\rho\) to be
trivial, so this conditional lemma has no additional
nontrivial torsion examples. Its proof above is given
independently precisely to preserve the logical
dependency used to establish R1.

\statement{E.51}{Proved}{The normalized section and the independent Magnus lemma remove the citation cycle}
The pure peripheral outer group \(\mathcal P\)
has a unique normalized homomorphic section into
\(\operatorname{Aut}(F)\) of the form
\[
 x\longmapsto x^\lambda,
 \qquad y\longmapsto f^{-1}y^\lambda f,
 \qquad f\in F'.
\]
For every prime \(p\), a finite-order continuous
automorphism of a finitely generated free pro-\(p\)
group which is identity on integral abelianization
is identity. This lemma is independent of LM4
D.25's final conclusion and of LM4 D.32.

\begin{proof}
For the first assertion, a peripheral outer class
has common multiplier \(\lambda\), since the
abelian relation \([x]+[y]+[z]=0\) equates the
three exponents. After an inner adjustment its
representative sends \(x\) to \(x^\lambda\) and
\(y\) to \(k^{-1}y^\lambda k\). Write
\([k]=(a,b)\in\widehat{\mathbb Z}^2\).
Conjugating that representative by \(x^a\) and
replacing its \(y\)-conjugator on the left by
\(y^{-b}\) gives
\(f=y^{-b}kx^{-a}\in F'\).
If two normalized representatives are inner-equivalent,
the inner conjugator centralizes \(x^\lambda\),
so is \(x^v\) by E.52. Their \(y\)-conjugators
then differ by left multiplication by \(y^w\)
and right multiplication by \(x^{-v}\).
The two zero abelian coordinates force \(v=w=0\).
Composition of normalized representatives has
conjugator \(f\alpha(g)\in F'\), proving the
homomorphism assertion as well as uniqueness.

For the second assertion, let the free basis be
\(x_1,\ldots,x_d\). Its integral Magnus map is
the continuous injective homomorphism into
\[
 \mathbb Z_p\langle\!\langle X_1,\ldots,X_d
               \rangle\!\rangle^{\times},
 \qquad x_i\longmapsto1+X_i.
\]
Continuity uses the \((p,X_1,\ldots,X_d)\)-adic
topology; injectivity follows already after reduction
modulo \(p\) from the faithful free pro-\(p\)
Magnus map. An automorphism identity on integral
abelianization substitutes
\(X_i\mapsto X_i+R_i\), with every \(R_i\)
of homogeneous degree at least two.
If some \(R_i\ne0\), choose the least occurring
degree \(m\). Induction on \(n\) shows that the
degree-\(m\) correction in the \(n\)-th iterate
is exactly \(n\) times the original correction:
any new substitution into a degree-\(m\) term
changes only larger degrees. The homogeneous
coefficient module is torsion-free over
\(\mathbb Z_p\), so no positive iterate can be
identity. This contradiction proves the lemma,
including \(p=2\).
\end{proof}

The acyclic dependency is now: E.52 supplies the
free-factor centralizers; E.51 supplies normalization
and the Magnus torsion lemma; E.50 supplies the
prime-order section calculation; the separate
three-cusp finite-cover argument in refereed
LM4 D.32 proves R1. Only after that may LM4 D.25
deduce triviality of its full multiplier-one
finite-order solutions. Its final R1-dependent
sentence must not be cited as an input to D.32.

\statement{E.52}{Known}{The precise Herfort--Ribes citation for the free peripheral reduction}
For \(F=\widehat F(x,y)\), each peripheral cyclic
group \(L=\overline{\langle u\rangle}\), where
\(u=x,y,(xy)^{-1}\), is malnormal, and
\[
 N_F(L)=L,
 \qquad C_F(u^a)=L
       \quad(0\ne a\in\widehat{\mathbb Z}).
\]
The analogous statements hold in \(F^{(p)}\),
with \(0\ne a\in\mathbb Z_p\). The precise
source in the free-factor reduction of LM4 A.7
is Herfort--Ribes (1985), \emph{Theorem B'},
Section~4, printed page~159.\cite{A_HR1985}

\begin{proof}
Theorem B' states that a factor of a free
pro-\(\mathcal C\) product intersects each of
its conjugates outside itself trivially, and
identifies the centralizer of a nonidentity
factor element with its centralizer in that
factor. Use the full formation of finite
groups, or the formation of finite \(p\)-groups.
Each of the three displayed \(u\)'s is a free
basis element: use \((x,y)\), \((y,x)\), or
\((x,(xy)^{-1})\), respectively. Hence
\(F=L\amalg L'\), and Theorem B' applies.
The power \(u^a\) is nonidentity because
\(L\simeq\widehat{\mathbb Z}\), or
\(\mathbb Z_p\). A normalizing element has
nontrivial factor intersection and therefore
belongs to \(L\). The reverse inclusions
follow from procyclicity.
\end{proof}

This citation correction concerns the free-factor
argument, which covers the pants group used in
LM4 D.28 and D.32. For a once-punctured surface
the boundary word is generally not primitive;
the separate peripheral commensurator theorem
used in LM4 A.7 cannot be replaced by this
free-factor citation. The locations in the
manuscript remain \S13 bis, folio~57 verso,
scan~118, and \S35, folios~253--254,
scans~499--501; E.50 also concerns \S31,
folios~198--199 verso, scans~391--394.

\bibliographystyle{apalike}
\bibliography{LM4_SUCCESSOR_references}
\end{document}
